%% LyX 2.3.6.1 created this file.  For more info, see http://www.lyx.org/.
%% Do not edit unless you really know what you are doing.
\documentclass[11pt,english,usenames,dvipsnames]{article}
\usepackage{mathpazo}
\usepackage[T1]{fontenc}
\usepackage[latin9]{inputenc}
\usepackage{geometry}
\geometry{verbose,tmargin=2.54cm,bmargin=2.54cm,lmargin=2.54cm,rmargin=2.54cm}
\setlength{\parindent}{2em}
\usepackage{color}
\usepackage{babel}
\usepackage{float}
\usepackage{amsmath}
\usepackage{amssymb}
\usepackage{graphicx}
\usepackage{setspace}
\usepackage[authoryear]{natbib}
\onehalfspacing
\usepackage[unicode=true,
 bookmarks=false,
 breaklinks=false,pdfborder={0 0 1},backref=false,colorlinks=true]
 {hyperref}
\hypersetup{
 hypertexnames=false,allcolors=OliveGreen}

\makeatletter
%%%%%%%%%%%%%%%%%%%%%%%%%%%%%% User specified LaTeX commands.
\usepackage{babel}
\usepackage{bbm}
\usepackage{placeins}
%\usepackage[osf]{mathpazo}

%\usepackage[usenames,dvipsnames]{xcolor}

\usepackage{fancyhdr}
\usepackage{etoolbox}
\usepackage[affil-it]{authblk}
\usepackage{array}
\newcolumntype{L}[1]{>{\raggedright\let\newline\\\arraybackslash\hspace{0pt}}m{#1}}
\newcolumntype{C}[1]{>{\centering\let\newline\\\arraybackslash\hspace{0pt}}m{#1}}
\newcolumntype{R}[1]{>{\raggedleft\let\newline\\\arraybackslash\hspace{0pt}}m{#1}}


\patchcmd{\@fancyhead}{\rlap}{\color{OliveGreen}\rlap}{}{}
\patchcmd{\headrule}{\hrule}{\color{OliveGreen}\hrule}{}{}
\patchcmd{\@fancyfoot}{\rlap}{\color{OliveGreen}\rlap}{}{}
\patchcmd{\footrule}{\hrule}{\color{OliveGreen}\hrule}{}{}

\usepackage{titlesec}

\titleformat{\section}
{\color{Black}\normalfont\Large\bfseries}
{\color{Black}\thesection}{1em}{}

\newcommand{\jtable}[4]{
\input{#1}
\hline\hline
\end{tabular}
\\ \vspace{0.1cm} \parbox{#4}{\footnotesize #3} 
\label{#2}
\end{table}
}

%\renewcommand{\footnotesize}{\fontsize{10pt}{11pt}\selectfont}

\makeatother

   \def\changesfont{\fontsize{5}{5}\selectfont}%
   \setlength{\marginparwidth}{20mm}
   \setlength{\marginparsep}{4pt}
      \usepackage{soul}
      \sethlcolor{cyan!30}%\hl{Text}
   \usepackage[commentmarkup=todo,todonotes={textsize=changesfont}]{changes}
   \definechangesauthor[name=CE, color=cyan]{CE}


\begin{document}
\title{\vspace{-1.5cm}Legal Reforms, Conditional Cash Transfers, and Intimate
Partner Violence: Evidence from Mexico\thanks{{\scriptsize{}This paper supersedes an earlier working paper titled
\textquotedblleft Conditional Cash Transfers for Women and Spousal
Violence: Evidence of the Long-Term Relationship from the Oportunidades
Program in Rural Mexico\textquotedblright. We thank Ruizhi Zhu for
excellent research assistance. We are grateful to Siwan Anderson,
Trinidad Beleche, Dwayne Benjamin, Matthieu Chemin, Sonia Laszlo,
Ben Nyblade, Torsten Persson, and especially Laura Chioda, Fabian
Lange, and Norbert Schady whose suggestions greatly improved the paper.
We also thank seminar participants at McGill University,
the Canadian Institute for Advanced Research IOG Program, the World
Bank, and the Inter-American Development Bank for helpful comments;
Ted Miguel and his research team for sharing their crime data. We
also gratefully acknowledge funding from the Canadian Institute for
Advanced Research, the Canada Research Chairs Program, the Inter-American
Development Bank, and the World Bank. The views expressed here are
those of the authors and do not necessarily reflect the views of the
funding institutions. We are responsible for any errors that remain.}}\\ \textbf{Word Count: 7300}}
\author{Gustavo J. Bobonis\thanks{{\scriptsize{}Professor, Department of Economics, University of Toronto,
BREAD, and J-PAL. Address: 150 St. George St., Toronto, Ontario, M5S
3G7, Canada. Tel: 416-946-5299. Fax: 416-978-6713. E-mail: gustavo.bobonis@utoronto.ca}} \hspace{1cm}Roberto Castro\thanks{{\scriptsize{}Professor, Regional Center for Multidisciplinary Research,
National Autonomous University of Mexico, Av. Universidad s/n, cto.
2, Col. Chamilpa, Cuernavaca, Morelos, 62210, México. Tel: (52-777)-329-1853.
E-mail: rcastro@crim.unam.mx}} \hspace{1cm}Juan S. Morales\thanks{{\scriptsize{}Associate Professor, Department of Economics, Lazaridis
School of Business and Economics, Wilfrid Laurier University. Address:
75 University Avenue West, Waterloo, ON, N2L 3C5, Canada. E-mail:
jmorales@wlu.ca}} \vspace{-0.3cm}}
\maketitle
\begin{abstract}
\onehalfspacing\textcolor{black}{\footnotesize{} We study the relationship
between divorce law reforms codifying intimate partner violence (IPV)
as legal grounds for unilateral divorce, the Oportunidades conditional
cash transfer program, and the incidence of IPV in Mexico. Using data
from three nationally representative surveys in 2003, 2006, and 2011,
we show the legal reforms lead to a 55 percent increase in annual
divorce rates, concentrated among couples with a history of violence.
Comparing groups of beneficiary and non-beneficiary households within
villages, we find that IPV rates converge for these couples in the
longer-run. Marital selection plays an important role in explaining
the \highlight[id=CE, comment={long-term ? or long-run ?\break ~\break This will continue, as US spelling is long-run more appropriate?}]{long-run} relationships.}{\footnotesize\par}

\noindent \textbf{\textcolor{black}{\footnotesize{}JEL Codes:}}\textcolor{white}{\footnotesize{}
}\textcolor{black}{\footnotesize{}J12, J16, K42}\textcolor{white}{\footnotesize{},
}\\
\textbf{\textcolor{black}{\footnotesize{}Keywords:}}\textcolor{black}{\footnotesize{}
divorce laws; conditional cash transfer programs; Oportunidades; divorce;
intimate partner violence; marital selection}{\footnotesize\par}
\end{abstract}
\pagebreak\doublespacing

\section{Introduction}

Intimate partner violence (IPV) is the most common form of violence
experienced by women globally: approximately 30 percent of women experience
this in some form during their lifetimes. IPV has been condemned internationally
as a serious and costly issue in human rights, public health, and women\textquoteright s
personal security (\citealt{devries2013intimate,fearon2014benefits}).
As a result, a growing number of national and subnational governments
have instituted multiple efforts to promote women's ability to exit
violent relationships as well as their empowerment within the household,
partly with the intent of curbing intimate partner violence and other
forms of violence against women.

Governments in many (mostly Western) countries have over the past
50 years introduced legal reforms that allow for unilateral divorce,
which ease individuals' ability to exit marital relationships.\footnote{These unilateral divorce law reforms, which allow a spouse to file
for divorce unilaterally and without the other spouse having committed
fault have been instituted across a number of states in the United States, in Canada,
and in European countries since the late 1960s.} Strategies more commonly adopted across the developing world have
focused on the implementation of legal reforms and social programs
that aim to enhance women's options outside of existing marriages
(i.e., property rights and inheritance reforms that favor women; cash
and in-kind transfers-based poverty alleviation programs targeted
to adult women), based on a growing consensus that targeting resources
to women can promote their empowerment within the household.\footnote{For a survey of the broad literature on policies aimed at promoting
female empowerment, see \citet{duflo2012women}. \citet{buller2018mixed}
and \citet{baranov2020theoretical} provide recent reviews of the
work on cash transfers and intimate partner violence in low and middle-income
countries. A growing body of research also studies how legal and judicial
systems as well as access to justice programs can address gender inequities
and improve women\textquoteright s rights and wellbeing, including
exposure to intimate partner violence (see, for example, \citealp{anderson2018legal,miller2019female,amaral2019gender,kavanaugh2019women}).
Surveys of this small but growing literature can be found in \citet{heise2011works,doyle2018economics}.} A literature evaluating this broad set of programs and policy reforms
in developing country settings tends to find moderate or substantial
short-term reductions in the incidence of IPV as a result of both
types of policy reforms \citep{buller2018mixed}. However, there is
limited evidence on the important issue of how local legal institutions
interact with social support programs to curb the incidence of IPV
\citep{baranov2020theoretical}.

The present paper investigates a series of questions regarding the
long-run relationship between divorce legislation, conditional cash
transfer programs and the incidence of IPV. Specifically, do legal
reforms that help ease individuals' ability to exit violent relationships
achieve \emph{sustained }reductions in the incidence of IPV in developing
country settings? If so, does this occur by inducing the marital dissolution
of violent relationships? Do conditional cash transfer programs complement
or substitute these legal efforts to curb intimate partner violence?

We study these questions in the context of rural Mexico. Starting
in 1997, state governments across the country adopted a series of
legal reforms formally including intimate partner violence as legal
grounds for unilateral divorce. Concurrently, the federal government
introduced in 1997 the landmark PROGRESA/Oportunidades/Prospera conditional
cash transfer program, one of the earliest and most comprehensive
poverty alleviation programs in Latin America. We exploit the introduction
of these policies to provide evidence of the short- and \highlight[id=CE, comment={short- and long-term. This style persists throughout. Consider changing all uses to short- and long-term ?}]{long-run} relationship
between these legal reforms, the Oportunidades CCT program, and the
incidence of intimate partner violence. We use data from three nationally
representative surveys: the National Surveys on the Dynamics of Household
Relationships (ENDIREH) of 2003, 2006, and 2011. These surveys collect
detailed information on IPV as well as households' program participation
and women's marital histories, and thus provide rich
data to document the consequences of the reforms for marital dissolution
and patterns of IPV across program beneficiary and non-beneficiary
households.

Our study uncovers several new findings regarding intimate partner
violence in Mexico and its relationship with instituted policies.
The legal reforms, which reduce the cost of exiting violent relationships,
lead to a substantial increase in divorce rates in the Mexican countryside.
Exploiting the staggered adoption of these legal changes across states,
we estimate a 0.35 percentage point (55 percent) increase in annual
divorce rates following the passage of the reform. \hl{These effects are
concentrated among Oportunidades program beneficiary households, those
in which women reporting having experienced intimate partner violence,
and those having experienced violence in their} \highlight[id=CE, comment={check structure of clause.}]{childhood}, as an important
predictor of IPV incidence. These patterns suggest that divorce and
marital selection can be an important driving force towards reductions
in the incidence of IPV.

Second, we study potential changes over time in the relationship between
receipt of the Oportunidades CCT program and IPV, across the 2003,
2006, and 2011 surveys. We follow the observational analysis in \citet{bobonis2013public},
comparing the incidence of women's experience of IPV across beneficiary
and non-beneficiary households within villages (in 2003), and construct
samples of comparable households for each of the two subsequent survey
rounds (2006 and 2011). We initially show that the short-run relationship
between receipt of the Oportunidades CCT and reductions in IPV documented
in \citet{bobonis2013public} are largely concentrated in states that
had adopted the legal reforms to the divorce laws before 2003, of which often are early
adoption states. The estimates suggest that beneficiary women in these
states were 18 percentage points less likely to be victims of physical
or sexual violence compared to non-beneficiary women in non-reform
states in the \highlight[id=CE, comment={short-run}]{short-run}. This finding suggests a strong \highlight[id=CE, comment={short-run}]{short-run}
complementarity between the two policies in reducing domestic violence.
In contrast with these \highlight[id=CE, comment={short-run}]{short-run} outcomes, we find no differences
in the incidence of IPV between women in beneficiary and non-beneficiary
households in the \highlight[id=CE, comment={longer-run}]{longer-run}. In combination with our results regarding
marital dissolution patterns, the evidence is indicative of the idea
that marital selection is an important mechanism of overall reductions
in IPV incidence over time.

Finally, in the appendix we present a simple model of intra-household
bargaining and conflict to help formalize the mechanisms via which
both reductions in the cost of divorce and income support programs
can affect equilibrium levels of marital dissolution and conflict
within the household. Building on the work of \citet{anderson2015suicide}
and \citet{brassiolo2016domestic}, our theoretical framework incorporates
bargaining and conflict with incomplete information, together with
endogenous divorce decisions, to formalize the effects that these
two distinct sets of policies can have on the incidence of conflict
and divorce. The first prediction we draw from the model shows how
legal reforms, through a reduction in the costs of divorce, increase
separations. More importantly, the model highlights the specific selection
mechanisms through which the implemented policies, both CCTs and legal
reforms, affect the $types$ of couples which remain in union. In
particular, the model reveals that relationships in which the private
gains from marriage are lower and the cost of conflict experienced
by the woman is higher, are more likely to dissolve once these policies
are introduced. If these characteristics are associated with the occurrence
of IPV, then these policies can increase the separation rate of relationships
in which violence is prevalent.

\subsection*{Relevant Literature}

Our paper contributes to the economics and policy literature on intimate
partner violence in a number of ways. Most importantly, we explore
the relationship between two separate factors, which have been documented
to affect IPV, legal frameworks and cash transfers. We highlight that
exploring the intersection of these mostly separate literatures can
help us better understand the effects of individual policies, and
in particular in understanding their long run implications.

A growing body of work has shown that unilateral divorce laws\textemdash{}instituted mainly in the context of developed countries\textemdash{}can help curb various forms of intimate partner violence (\citealt{dee2003until,stevenson2006bargaining,brassiolo2016domestic}).
However, it is unclear whether the evidence provided by these studies
is generalizable to developing countries, where differences in both
socio-economic and health conditions as well as in cultural views
towards IPV and social norms regarding gender roles (especially those
that link notions of manhood to dominance and aggression) may limit
the effectiveness of these legal rules in improving women's status
and conditions within the household. \citet{GARCIARAMOS2021102623}
explores the impact of unilateral divorce laws in Mexico and finds
increased physical, economic and emotional violence in the long-run
for women who remained married, as well as a positive relationship
between the adoption of unilateral divorce laws and divorce rates.
Similarly, both \citet{aguirre2019changes} and \citet{hoehn2021impact}
find that unilateral divorce laws significantly increased divorce
rates across Mexico between 2008 and 2017. Our study complements this
work: it examines a different set of (earlier) legal reforms\textemdash{}instituted in most states in the country\textemdash{}that also enhance
women's rights by creating a more equitable legal environment for
marriage and divorce, while exploring how these reforms interacted
with the Oportunidades CCT program. \citet{beleche2017domestic} studies
the consequences of these reforms for extreme measures of intrahousehold
violence, particularly female and male suicide rates,
and finds no evidence of reductions in these extreme forms of violence
and conflict resolution. Using aggregate administrative data on divorce
rates, \citet{beleche2011til} find no overall effect on divorce rates due to this legal
reform.\footnote{In contrast, we document increased divorce rates in the specific sub-population
we study (young women in rural settings), detailed below. }

A parallel and significant body of work has studied the effects of
conditional and unconditional cash transfer (CT) programs on intimate
partner \highlight[id=CE, comment={insert reference to study here, at end of sentence.}]{violence.} While this body of work finds that some subgroups are at increased risk of some forms of violence, the collective evidence
supports the view that CTs lead to short-run reductions in IPV \citep{buller2018mixed,baranov2020theoretical,heath2020cash}.
Few studies document the long-run consequences of the introduction
of CT programs. \citet{perova2010buying}, \citet{ritter2014te},
and \citet{diaz2022unconditional} study the consequences of the introduction
of the Juntos CT program in Peru; all find significant reductions
in the incidence of physical and sexual violence over a period of
up to five years following the introduction of the program. \citet{roy2019can,roy2019transfers}
study whether cash transfer programs lead to persistent reductions
in IPV once these programs end. They show evidence based on a randomized
controlled trial in rural Bangladesh that cash or food transfers to
poor women alone do not lead to sustained reductions in IPV following
the end of the \protect\highlight[id=CE, comment={clarify author names in footnote as well as in text}]{program.}\hl{\footnote{\hl{They} presented strong evidence that a complementary nutrition behavior-change communication intervention, together with cash transfers in
particular, led to sustained reductions in IPV four years after the
end of the program. Their evidence suggests that the combined intervention
led to more sustained impacts through persistent increases in women\textquoteright s
bargaining power, men\textquoteright s costs of perpetrating violence,
and poverty-related emotional well-being.}} Relatedly, \citet{litwin2019conditional} document that transfers
from the Bolsa Familia program to Brazilian municipalities are associated
with separations of couples (and especially separations of couples
with children), highlighting an important mechanism through which
CTs may reduce \protect\highlight[id=CE, comment={clarify author names in footnote as well as in text}]{IPV.}\hl{\footnote{However, \hl{they} find no relationship between the CT program and female
homicides (the most extreme form of IPV).}} 

\hl{Although previous literature has mostly highlighted how policy efforts
affect the behaviour of violent partners, the potential of these to
reduce IPV by inducing divorce remains a relatively unexplored channel
through which this objective could be} \highlight[id=CE, comment={these studies - this objective : it is unclear what these and this are referring to here, consider expanding clause for clarity.}]{achieved.} Our findings suggest
that the short-run causal effect of CTs on IPV, documented by the
existing literature, may not necessarily persist over time if changes
in the legal environment affect household formation and dissolution
patterns in such a way as to independently curb intimate partner violence.
Ours is the first study to explore the extent to which legal reforms
and social policies, which may both have independent effects on IPV,
work as complements or substitutes. The results are consistent with
the view that policies to promote the improvement of women's livelihoods
within the household can lead to sizeable reductions in IPV in the
\highlight[id=CE, comment={short-run / short-term}]{short-run}, but that marital selection can be a significant mechanism
for explaining such effects in the \highlight[id=CE, comment={longer-run / change to longer run. use hyphen for short-term and no hyphen for short run}]{longer-run}. The evidence we present
suggests that the policies can be important complements in reducing
IPV.

\section{Background}

\subsection*{Family Violence Divorce Law Reforms}

In Mexico, the laws regarding marriage, divorce, and family affairs
are embedded in state civil codes; each of the 32 states has its own
civil code regulating these proceedings. State civil codes have historically
stipulated motives such as adultery, and out-of-wedlock births (among
others) in addition to mutual consent, as grounds for divorce. In
the late 1990s, states began adopting legal reforms with the objective
of modernizing civil codes, one of which included adding family violence
as a cause for divorce. The precise wording of the civil code varies
across states, but it generally allows partners to unilaterally initiate
divorce proceedings claiming family violence\textemdash{}the committal
of violence by one of the partners against the other, their mutual
children, or the children of one of them\textemdash{}as the cause.
The specific dates at which these legal reforms were adopted across
Mexican states are listed in Figure 1 (the source is from the coding carried
out by \citealp{beleche2017domestic}). The first state to adopt the
reform was the Federal District (Mexico City) in December 1997.
By 2003, the majority of states had adopted the legal reform. In that
year, domestic violence or threats as the cause of divorce accounted
for about 1.6 percent of all divorces in urban areas, and about 2.2
percent in rural areas.\footnote{Authors' own calculations based on data from INEGI: http://www.beta.inegi.org.mx/programas/nupcialidad/.}

Figure 2 shows trends in the adoption of family violence reforms
and divorce rates since 1997, suggesting a relationship between
the increase in divorce rate and the divorce law reform. In 1997 less
than 10 percent of the states allowed domestic violence as grounds
for divorce, but by 2001 this proportion had increased to almost 50
percent, and to almost 70 percent by the year 2006. Coinciding with
the passage of divorce legislation, divorce rates also increased.
In 1993 there were 0.41 divorces per thousand persons. This statistic
rose to 0.76 divorces per thousand persons in 2005. According to Mexican
statistical yearbook statistics, in 2001\textendash{}2006 mutual consent accounted
for over 70 percent of divorces, while separation or abandonment accounted
for 5\textendash{}10 percent. In addition, while the majority of divorces are
judicial, there has also been a rise in the number of administrative
divorce filings over time. Relative to the other causes for divorce,
the proportion of divorce filings listing domestic violence as the
cause is slightly over one percent. IPV, however, is a major public
health concern affecting Mexico. Based on household surveys in 2003,
approximately 44 percent of women living with a partner reported
having been a victim of domestic violence (ENDIREH, 2003). When intimate
partner violence is allowed as grounds for unilateral divorce, an
abused spouse may threaten to use it against the other in order to
obtain mutual consent to dissolve the marriage.

\subsection*{Oportunidades Program}

The Mexican government initiated in 1997 a conditional cash transfer
program named PROGRESA, renamed \textit{Oportunidades} in 2001 under
the Fox Administration (\hl{and renamed PROSPERA under the Pe$\~n$a Nieto}
\highlight[id=CE, comment={"renamed PROSPERA in (insert date) under the Pea Nieto administration" ?}]{administration}), aimed at alleviating poverty and improving the human
development of children in rural Mexico. The program targets the poor
in marginal communities, where 40 percent of children from poor
households drop out of school after the primary level. The program
has expanded considerably since its inception and has become an integral
component of Mexico's social development and poverty reduction efforts.
As of 2013, Oportunidades provided cash transfers to 6.5 million families,
conditional on children school attendance, health checks, and participation
in health clinics.

The targeting of the program was done at two levels. First, eligible
localities were identified on the basis of a locality-level eligibility
rule. Program officials used locality-level characteristics from the
Mexican 1995 Mini-Census of Population to construct a marginality
index for each locality that reflected its degree of marginalization
and was correlated with the community's incidence of poverty.\footnote{The variables used to construct this marginality index were: (i) the
locality's population, (ii) the number of dwellings in the village,
(iii) the proportion of the adult population who was illiterate, (iv)
the proportion of adults working in the agricultural sector (in 1990),
the proportion of households (v) without potable water, (vi) without
drainage, (vii) without electricity, (viii) with a dirt floor (in
1990), and (ix) the average number of persons per room in each household
(in 1990).} Second, program enumerators conducted household surveys within eligible
localities to identify households that would be classified as poor.
Based on asset holdings used as proxy variables for poverty, the program
administrators generated a proxy-means \highlight[id=CE, comment={In footnote 8, change to "of poor and" or "classification as poor and non-poor" ?}]{test.}\hl{\footnote{Within a sub-sample of communities, a poverty indicator was constructed
using household income data collected from baseline surveys. A discriminant
analysis was then separately applied in each region in order to identify
the household characteristics that maximized the correct classification
\hl{of as poor} and non-poor (minimizing Type I and Type II targeting errors).
Eligible households were identified on the basis of this welfare index
(see \citealt{skoufias2001targeting} for a more detailed description
of the targeting process).}} Therefore, within each eligible community, only households below
a threshold became program beneficiaries. The list of potential beneficiaries
was then discussed in a community meeting and suggested revisions subsequently were
sent to the central Oportunidades office. In practice, very few changes
were made to the list of targeted households \citep{skoufias2001targeting}.
This targeting and program eligibility information is an important factor in
the construction of our sample of eligible women (see Sections IIII
and IV.B).

Initially, a locality was eligible for Oportunidades if it was classified
as \textquotedbl poor\textquotedbl{} (marginality grade 4) or \textquotedbl very
poor\textquotedbl{} (marginality grade 5) out of a 1\textendash{}5 scale based
on the locality-level marginality index, and if it had access to a
primary school, a secondary school, a health center, and was classified
as rural (defined as inhabited by fewer than 2,500 people), but had
at least 50 inhabitants \citep{skoufias2001targeting}. The last criterion
was relaxed early on to incorporate some semi-urban localities (localities
with between 2,500 and 14,999 inhabitants). The health center criterion
was relaxed in 1998 when mobile health clinics were introduced. Since
then, the inclusion of less marginal localities into the program has
been gradually extended. Between 2000 and 2011, the program's coverage
expanded from around 53,000 localities and 2.5 million families, to
97,000 localities and 5.8 million families. The program was phased-in
through a different targeting design in urban areas starting in 2001.
Since this targeting mechanism is very complex and substantially different
to the one implemented in rural and semi-urban areas, and in order
to maintain a sample comparable to that of the short-run study, we
focus our analysis on rural households. We present additional information
on the Oportunidades program in the appendix.

\section{Data, Measurement, and Summary Statistics}

\subsection*{Description of the ENDIREH Surveys and Study Samples}

We use data from Mexico's National Surveys on Relationships within
the Household (\textquotedbl Encuesta Nacional sobre la Dinámica
de las Relaciones en los Hogares\textquotedbl, or ENDIREH) of 2003,
2006 and 2011. These are \hl{three nationally representative cross-sectional household }\highlight[id=CE, comment={note my change here..}]{surveys} measuring the prevalence and intensity of intimate
partner violence, among other intra-household interactions. They contain
data on household demographics, socio-economic characteristics, (limited)
marital histories, household decision-making, marital conflict, and
a module designed to measure the prevalence and severity of intimate
partner violence. The 2003 survey was administered to 54,230 women
15 years or older living with a husband or partner, whereas the 2006
and 2011 surveys were administered respectively to 113,561 and 152,636
women in the same age range but independent of marital status.

We construct measures of the incidence of violence that consist of dichotomous
variables indicating whether the female partner had suffered physical
or sexual violence from her spouse or partner in the past 12 months.\footnote{These definitions follow closely the documentation and results of
the survey in \citet{castro2006violencia}. A previous version of
this paper also included emotional abuse as an outcome, but analogously
with the results for physical and sexual violence, we see no differences
in the incidence of emotional violence between beneficiary and nonbeneficiary
women in the 2006 and 2011 surveys. These results are available upon
request.} Physical violence includes pushing, kicking, throwing objects, hitting
with hands or objects, choking, attacking with a knife or blade, and
shooting. Sexual violence includes demanding sex against the woman's
will, forced sexual acts, and forced sexual relations. A single incident
reported within the past year is classified as violence.

Data on program participation come from the ENDIREH surveys and are
self-reported by women. The measure of program participation available
in the ENDIREH 2003 is based upon whether the woman receives benefits from any
government support program at the time of the survey. Although Oportunidades
is the largest and most generous cash transfer program, there are
other small government programs that provide non-cash benefits. As
a result, this measure may over-report the receipt of Oportunidades
benefits. However, although there is some noise in the data (since
only ten households per village are randomly selected to participate
in the survey) the correlation of the proportion of beneficiary households
using the ENDIREH survey data with administrative data on the number
of recipient households at the locality level in 2003 is 0.84 (not
reported in the tables), which suggests that the information from
the household survey closely represents receipt of Oportunidades benefits.
In addition, since program receipt is measured at the time of the
survey, it combines couples who have received the transfer for varying
lengths of time (which we unfortunately do not observe) and our results
should be interpreted with this in mind.

The ENDIREH 2006 and 2011 surveys ask women specifically whether they
receive benefits from Oportunidades, and separately whether they are
beneficiaries of other government support programs. In order for the
analysis to be comparable to that using data from the ENDIREH 2003
survey, the measure we use is the analogous measure of being a beneficiary
from any government support program (i.e., Oportunidades or other).
In the samples for the ENDIREH 2006 and 2011 selected for our analysis,
only between 1.6 and 3.0 percent of those who report being beneficiaries
of any government support program, report not being beneficiaries
of Oportunidades. These reliability checks suggest that the information
from the household survey closely represents receipt of Oportunidades
benefits.\footnote{We also estimate models using the ENDIREH 2006 and 2011 data with
the Oportunidades beneficiary indicator as the explanatory/treatment
variable of interest. The results do not differ in any significant
way from those reported in the tables. These are available from the
authors upon request.}

In order to minimize potential selection bias as a result of the
targeting and endogenous take-up of the program, we restrict the analysis
of the short-term relationship in 2003 to a particular subset of households
as in \citet{bobonis2013public}. The 2003 sample includes couples
with women 25 years or older, with children younger than 11 years
old, and who have been married since at least 1997. These restrictions
result in a sample of 2,613 couples. For our analysis in subsequent
years, we first construct a \textit{pseudo-panel} of comparable households.
We restrict the 2006 (2011) survey sample to couples with women 28
(33) years or older with children between the ages of 3 and 13 (8
and 18) years, who have been together since 1997. The resulting overall
sample sizes for the pseudo-panel are 4,240 in the 2006 survey and
5,208 couples in the 2011 survey. As we will discuss below, these
sample restrictions minimize potential confounding due to the endogenous
take-up of the program based on household socio-economic characteristics
and preferences for human capital investments (see Section 4).

In an alternative and complementary exercise, we construct a \textit{replication
sample }for the subsequent surveys that consists of women with the
same characteristics at the time of the survey, that is: women 25
years or older, with children younger than 11 years old, who have
been in a relationship for at least six years (this last restriction
is analogous to women being married since 1997 for the 2003 sample).
Note that while the pseudo-panel approach tries to maximize the overlap
of women across the samples, this alternative approach will include
many new women and will exclude others that no longer meet the selection
criteria (for instance, if the children are now out of primary school).
The resulting sample sizes for the replication sample are 4,318 for
the 2006 survey and 5,931 for the 2011 survey.

It is worth highlighting that each of these samples suffer from specific
limitations. In particular, the relationship between receipt of the
CCT and IPV estimated using the pseudo-panel conflates the policy
impacts with cohort effects and the time path of divorce, while that
obtained using the replication sample conflates the policy impacts
with changes in eligibility and the expansion of the program. However,
as we discuss below, the descriptive patterns we observe are very
similar in both subsamples.

Finally, to analyze the effect of the IPV divorce law reforms on the
divorce rate in our sample, we use the 2011 ENDIREH survey and construct
a \textit{state-level panel}. We start with the 2011 pseudo-panel,
the sample of women who remain in union in 2011 and who match our
sample restrictions (5,208 women), and pool it with the sample of
women who divorced between 1998 and 2011 but who share the same restricted
characteristics as the initial sample (in 2003 these women were 25
years or older and had children younger than 11, consisting of 472
women). The pooled sample has 5,680 observations. We use data on the
years of divorce in the pooled sample to estimate a divorce-rate for
each state-year cell. We do this for the overall sample, as well as
separately for both beneficiaries and non-beneficiaries of Oportunidades
(at the time of the survey). The divorce-rate panel consists of 448
observation cells (32 states x 14 years).

\subsection*{Descriptive Statistics}

\textcolor{black}{Spousal violence remains a pervasive phenomenon
in rural Mexico, but one that has decreased considerably throughout
the period (Table 1). Whereas 16.5 percent of women in the sample
reported experiencing some form of physical or sexual spousal violence
in the year 2003, the incidence had decreased to 14 percent by the
year 2006 and to 10.4 percent by 2011, in the pseudo-panel sample.
The reductions are even larger for the younger cohorts; in the replication
sample the incidence was 13.5 and 8.8 percent for 2006 and 2011, respectively.
The patterns in our selected sample are slightly more pronounced but
generally consistent with overall trends in Mexico.}\footnote{According to estimates based on the ENDIREH 2003-2011 survey rounds,
the overall incidence of physical violence among women was estimated
at 9.3 percent in the year 2003, increased slightly to 10.3 percent
in the year 2006, but decreased substantially to 4.4 percent in the
year 2011. The overall incidence of sexual violence has decreased
consistently across rounds\textemdash{}from 7.8 percent in the year
2003 to 6.0 and 2.8 percent respectively in the years 2006 and 2011
\citep[193]{casique2014expresiones}.}\textcolor{black}{{} }

\textcolor{black}{Households in the sample are of relatively low socio-economic
status. More importantly, we observe some stark differences in a number
of dimensions of socio-economic status as we compare the samples of
couples across survey years. A significant share of women report speaking
an indigenous language (14 percent in 2003, 17/16 percent in 2006,
and 21/19 percent in 2011, for the pseudo-panel/replication sample);
this ethnic identity is highly correlated with low socio-economic
status in Mexico (Table 2, Panel A). In addition, approximately nine
percent of women in 2003 have no schooling, and this figure increases
to 15 percent among the pseudo-panel couples selected in 2011. The
average age of women in the sample is 35.4 years in 2003, 38/35.2
years in 2006, and 43/35.1 years in 2011. The trend in age for the
pseudo-panel is explained by the age restrictions imposed on the samples
and is consistent with our attempt of following the same set of women
over time. Male partners belong to the same age group (the average
partner age is 38.4 years in 2003, 46.5/38.7 years in 2011), have
similar schooling attainment levels, and are as likely to have an
indigenous background (Table 2, Panel B). Households are relatively
large, with more than 5 members on average, a statistic usually correlated
with low socioeconomic status in the Mexican context.}

\textcolor{black}{The proportion of women who report having been exposed
to intra-household violence between their parents during childhood
is quite high, at approximately nine percent in 2003, 11/10 percent in
2006, and 13/11 percent in 2011 (Panel A). Given the existing concerns
and evidence regarding the intergenerational transmission of violent
behavior, this suggests that women in this context may be at a particularly
high risk of experiencing spousal violence, helping explain the prevalence
of violence reported above. The proportion of women reporting that
their male partners were exposed to spousal violence between their
parents during childhood is also significant, but decreases, from
18 percent in 2003, to 13/12 percent in 2006, and to 12/10 percent
in 2011. These are also important predictors of spousal abuse among
current partners (e.g., \citealt{bowlus2006domestic,casique2014expresiones}).}

\section{Empirical Methodology}

\subsubsection*{Relationship between Divorce Laws and Divorce Rates}

To study the relationship between the introduction of the IPV divorce
laws and the incidence of divorce we exploit the differential timing
at which these reforms were introduced across Mexican states in a
difference-in-differences framework. We use data on marital and family
histories from the last wave of the ENDIREH survey (2011) to identify
all women who matched our pseudo-panel characteristics (older than
25 and with young children in 2003, in union in 1997). With this information
we construct a state-year panel of divorce rates as described in Section
3. We then estimate regressions of the following form:

\begin{equation}
Y_{st}=\alpha+\theta PostDL_{st}+\gamma_{s}+\gamma_{t}+\varepsilon_{st}
\end{equation}

where $Y_{st}$ is the divorce rate in state $s$ and year $t$, and
$PostDL_{st}$ is an indicator equal to 1 if the divorce law had been
introduced in state $s$ by year $t$. Our preferred specification
includes state fixed effects $\gamma_{s}$ to control for time-invariant
heterogeneity across states and year fixed effects $\gamma_{t}$ to
control for time-shocks shared by all states. In some specifications
we also include linear state time trends.

In addition, we use variation across states on whether the reform
had been introduced in 2003 (year of the first survey wave) to investigate
whether the introduction of these reforms led to differential divorce
rates depending on respondent characteristics: violence in woman's
childhood, violence in marriage, and Oportunidades beneficiary status.
The baseline regression equation is the following:

\begin{equation}
Y_{is}=\alpha+\theta_{0}DL2003_{s}+\theta_{1}DL2003_{s}*V_{i}+\theta_{2}V_{i}+X_{is}\beta+\varepsilon_{is}
\end{equation}

where $Y_{is}$ is and indicator equal to 1 if individual $i$ from
state $s$ was divorced in the measured year (2003, 2006, or 2011);
$DL2003_{s}$ is equal to one if the IPV divorce law had been introduced
by 2003 in state $s$, $V_{i}$ is the covariate of interest, $X_{is}$
are the pre-determined covariates, and $\varepsilon_{is}$ are unobserved
determinants of divorce. Regressions are weighted by inverse sampling
weights aggregated at the state level. We cluster standard errors
at the state level.

\subsubsection*{Relationship between Divorce Laws and Intimate Partner Violence}

To investigate the relationship between spousal abuse and IPV divorce
laws, we estimate models analogous to those from equation (2) but
with IPV as our dependent variable. In addition, we estimate a similar
model which includes village fixed-effects that absorb time-invariant
characteristics at the village level. Because the village fixed effects
are linearly dependent with the IPV divorce law indicator, in this
case we can only identify the differential relationship between the
reforms with IPV depending on women's characteristics. Specifically,
we estimate the following equation:

\begin{equation}
Y_{isv}=\theta_{1}DL2003_{s}*V_{i}+\theta_{2}V_{i}+\alpha_{v}+X_{isv}\beta+\varepsilon_{isv}
\end{equation}

where $Y_{isv}$ is and indicator equal to 1 if individual $i$ from
village $v$ in state $s$ was a victim of physical or sexual violence;
$DL2003_{s}$ is equal to one if the IPV divorce law had been introduced
by 2003 in state $s$; $V_{i}$ is the covariate of interest, $\alpha_{v}$
are village fixed effects, $X_{isv}$ is the vector of pre-determined
covariates, and $\varepsilon_{isv}$ are unobserved determinants of
divorce. We do the analysis for each of the 2003 sample, and the pseudo-panel
and replication samples for 2006 and 2011. We cluster standard errors
at the state level.

\subsubsection*{Relationship between Oportunidades Beneficiary Status and Intimate
Partner Violence}

To obtain estimates of the relationship between Oportunidades beneficiary
status and the incidence of spousal abuse, we follow the analysis
in \citet{bobonis2013public} and use ordinary least squares models
conditioning on a large set of pre-determined individual and household
socio-economic characteristics, as well as locality/village fixed
effects, in order to capture any village-specific unobserved heterogeneity
influencing spousal abuse patterns (e.g., access to health clinics,
community groups, village-level conditions affecting partners' socio-economic
conditions and economic opportunities). The regression equation for
outcome $Y_{iv}$ is the following: 
\begin{equation}
Y_{iv}=\theta T_{iv}+X_{iv}\beta+\alpha_{v}+\varepsilon_{iv}
\end{equation}
where the treatment indicator $T_{iv}$ equals one for beneficiary
household $i$ in village $v$ and is zero otherwise; $X_{iv}$ are
the pre-determined covariates that are possibly significantly correlated
with $T_{iv}$ and $Y_{iv}$; $\alpha_{v}$ are village fixed effects,
and $\varepsilon_{iv}$ are unobserved determinants of domestic violence.
Since $T_{iv}$ is measured at the time of the survey, our estimates
combine possible short-term relationships, among couples receiving
program benefits for short periods, and those that have been receiving
program benefits for longer. We cluster standard errors at the village
level.

\subsubsection*{Dealing with Endogenous Selection into the Oportunidades Treatment}

We conduct the analysis using various strategies to minimize the extent
of selection bias in our estimates. First, as mentioned in \highlight[id=CE, comment={Section 3}]{Section
III}, our pseudo-panel uses a sub-sample of households with children
ages between 0 and 10 in 2003 (ages between 3 and 13 in 2006, and
ages 8 and 18 in 2011) and who have been in union since 1997. Second,
we condition on a set of pre-determined individual and household socio-economic
characteristics which are strongly correlated with determinants of
program eligibility and likely capture a large component of the variation
determining program take-up, thus minimizing concerns of endogenous
program take-up. In addition, we restrict the sample to women ages
25 and older in 2003 (28 and older in 2006, 33 and older in 2011).
Most importantly, to address the endogenous targeting of the program
to poor communities, we make comparisons of beneficiary and non-beneficiary
households within villages in order to remove all selection based
on the village-level targeting of the program. This within-village
comparison dramatically reduces the observed selection into the program
\citep{bobonis2013public}. In the appendix we discuss the potential
concerns of unobserved heterogeneity for the within-village household
comparison.

\section{Results}

\subsection*{Relationship between Family Violence Divorce Laws and Divorce Rates}

We begin our analysis by looking at the relationship between the legal
reforms instituted across Mexican states and divorce rates. The results
from estimating equation (1) are presented in Table 3. We present
estimates for the overall sample, as well as for specific subsamples
of women. Our preferred estimate includes both state linear time trends
and other legal reforms related to domestic violence which were instituted
around the same time.\footnote{These include criminalization of domestic violence and the establishment
of prevention and assistance programs for victims of domestic violence.
These policies are documented and studied in \citet{beleche2017domestic}.} The analysis suggests that legal reforms led to a statistically significant
increase in divorce rates across adopting states. The average annual
divorce rate in our sample is 0.64 per 100 couples. The estimated
coefficient suggests an increase of 0.37 divorces (column 2), or about
60 percent. Women who are beneficiaries of Oportunidades have slightly
lower divorce rates (0.59 per 100 couples), but experience a 0.49
rate increase following the reforms. Women who report violence in
their marriage have significantly higher divorce rates on average,
2.48 for 100 couples per year. We also estimate larger increases in
the divorce rates following the institution of the legal reform: a
1.39 divorce rate increase (significant at the 90 percent level),
a 56 percent increase in proportional terms.

In a complementary exercise we present estimates from an event-study
design in which we allow the coefficients to vary flexibly relative
to the timing of the reforms. We present these estimates in Figure
3, including an analysis in which we estimate divorce rates separately
for beneficiary and non-beneficiary women. The estimates reveal a
pattern consistent with those immediately above and of which suggest that the effects of
the legal reform on divorce was concentrated on beneficiary women.

\hl{The second part of our analysis presents estimates from the regression
model outlined in equation (2). This analysis using individual data exploits one particular threshold regarding the timing of the
legal reforms, as we compare states that had, by 2003 instituted the
divorce law reforms, relative to those which }\highlight[id=CE, comment={please note my changes here.}]{had not.} The advantage
of this framework is that it allows us to include our control variables
in the estimating equation. The results are presented in Table 4.
Women in states that had instituted the reforms by 2003 were more
likely to be divorced by 2006 (column 2, statistically significant
at the 90 percent level) and by 2011 (column 6, not statistically
significant), by about 1.3 percentage points. Interacting the presence
of the reform with the characteristics of these women suggests that the
relationship between the legal reform and the propensity to divorce
is heterogeneous across different subgroups of women. In particular,
women who experienced violence in their households as children (columns
3 and 7), women who experience violence in their marriage (columns
4 and 8), and women who are beneficiaries of Oportunidades (columns
5 and 9), are more likely to be divorced in legal reform states by
2006 and by 2011, relative to their converse groups in non-reform
states. The results from this complementary analysis confirm the findings
from state-panel regressions.

\subsection*{Relationship between Divorce Laws and Intimate Partner Violence}

The results from this analysis are presented in Table 5. Columns 1,
4, and 7 evaluate whether women who lived in states that had passed
the legal reforms by 2003 experienced spousal violence at different
rates relative to women in non-reform states. We find no significant
relationship between IPV divorce laws and spousal violence in this
cross-sectional analysis. This suggests that the reforms themselves
did not substantially change IPV rates and/or that states with higher
rates of IPV were not necessarily the first ones to pass these reforms.

We next test whether the relationship between IPV divorce laws and
spousal violence is heterogeneous depending on whether women experienced
violence in their homes as children. In 2003, women who were in reform
states and who had experienced violence as children were between 14
and 17 percentage points more likely to be victims of spousal abuse
(columns 2 and 3) relative to women who had not experienced violence
as children. We also observe that the relationship between violence
in childhood and domestic violence appears significant in the 2006
and 2011 samples but is not restricted to reform states.

\subsection*{Relationship Between Oportunidades Beneficiary Status and Intimate
Partner Violence}

We start the presentation of the relationship between Oportunidades
and IPV with a graphical analysis of patterns in the data. Figure
4 shows the trends in physical violence among couples across the three
survey years. The incidence of physical abuse among women in non-beneficiary
households is high at almost 13 percent in 2003, and presents a
downward trend over time to around 10 percent in 2006, and to around
7 percent in 2011 (significant at the 90 percent confidence level).
In comparison, the incidence among beneficiary couples is around 9
percent, lower than that among non-beneficiary couples in 2003 (4
percentage points, not significant). However, the incidence among
beneficiary couples hovers around 10.0 percent in 2006 and 7 percent
in 2011, such that physical violence rates among these two groups
of households converge in the longer-run. We observe a similar, although
less stark, pattern of short-run differences (in 2003) and later convergence
(in 2006 and 2011) in the incidence of sexual violence.

\textcolor{black}{Estimates of the overall five-year (2003), eight-year
(2006), and thirteen-year (2011) relationship between program beneficiary
status and spousal violence outcomes are displayed in Table 6. For the
purposes of comparison, we start the discussion with our preferred
estimates of the short-run relationship. As documented in \citet{bobonis2013public},
domestic violence incidence rates in the short-run are significantly
lower among beneficiary couples than among non-beneficiary couples (column
1), the estimated difference in the incidence of physical or sexual
abuse being 12.5 percentage points (75 percent).} \textcolor{black}{We
then investigate whether this relationship is heterogeneous between
states which had instituted the IPV divorce laws by 2003 and those
that had not (column 2). Although we do not have the precision to
detect a differential relationship (the interaction term between Oportunidades
beneficiary and IPV divorce law states is large but not statistically
significant), the overall relationship in IPV law states is statistically
significant and large, suggesting that women in reform states who
were beneficiaries of Oportunidades were 18.5 percentage points less
likely to be victims of physical or sexual abuse relative to non-beneficiary
women in reform states. The analysis suggests that there is a strong
short-run complementarity between the Oportunidades program and the
legal reform in terms of reductions of IPV, and that the short-run
reductions in domestic violence previously documented were concentrated
in these reform states.}

\textcolor{black}{In contrast to the short-run relationship, domestic
violence incidence rates do not vary significantly between beneficiary
and non-beneficiary households in 2006 and 2011 (columns 3\textendash{}6). Our
preferred estimate for 2011 shows a statistically insignificant difference
in the incidence of physical and sexual violence of 1 percentage point
(column 5).}\footnote{\textcolor{black}{We also find no evidence of a significant difference
in the incidence of violent threats or acts of emotional violence
among beneficiary women, a finding documented in a previous version
of this article \citep{bobonis2015conditionalc}.}}\textcolor{black}{{} The heterogeneity analysis for 2006 suggests that
beneficiary women in reform states were 4.1 percentage points }\textit{\textcolor{black}{more}}\textcolor{black}{{}
likely to be victims of physical or sexual abuse relative to non-beneficiaries
in non-reform states (column 4). The finding could be partially explained
by a combination of the strong selection mechanism combined with a
``male backlash'' effect. This relationship is, however, not present
in neither the replication sample for 2006 (Panel B), or either of
the 2011 samples (column 6). The interaction coefficient for 2011,
though statistically insignificant, is, as in 2003, negative. In some
of the specifications which include additional control variables,
this coefficient becomes statistically significant (shown in Table
8 and discussed in the appendix), suggesting that to some extent,
this complementarity between IPV divorce laws and the CCT program
persists through to 2011.}

\subsection*{Discussion, Repeated Interactions and Marital Selection}

\textcolor{black}{The stark differences in the longitudinal pattern
of the relationship suggests that the models of violence and household
bargaining, in which male partners may use violence as instruments
of coercion \citep[see][]{bloch2002terror,bobonis2013public,anderson2015suicide}
may correctly capture short-run interactions within the household
but may do so poorly in the longer-run. In this class of models male
partners are heterogeneous in their willingness to engage in violence
and have private information regarding the ``gains to
marriage'', such as their own private income or their
status within the household based on traditional gender roles. However,
once this information has been revealed through partners' actions,
couples with violent types may dissolve such that couples in future
periods are disproportionately composed of non-violent types. This
selection patterns can lead to both a tendency for violence rates
to decrease among couples remaining in union over time and for their
relationship with program receipt status to be dampened.}

\textcolor{black}{The evidence we present is consistent with this
interpretation. First, we have documented that the short-run relationship
with Oportunidades receipt is concentrated in reform states. We have
also shown that women in reform states experienced an increase in
marital dissolution rates over this time period. Furthermore, the
increase in marital dissolution rates was larger for beneficiary women
and women who experienced violence in their marriage. This increase
in divorce-based selection can help explain both the drop in the incidence
of violence among these cohorts of women and the absence of a relationship
between the Oportunidades program and the incidence of IPV in the
subsequent surveys. }In the appendix we present a formal model of
spousal relations that highlights the mechanisms through which legal
reforms and conditional cash transfers affect divorce and conflict
within the household, and reveals how different types of couples remain
in union when the policies are implemented.

\textcolor{black}{Second, a comparison of socio-economic and demographic
characteristics of households across the three survey waves suggests
that there are important changes in their distributions (see Table
2). Women are more likely to report speaking an indigenous language
in the later survey waves (which is correlated with low socio-economic
status in Mexico), they tend to have lower educational attainment
levels, and report a higher prevalence of violence in their households
during childhood. Moreover, their male partners are more likely to
be indigenous themselves. However, they report a lower prevalence
of violence in their male partners' households during childhood -
decreasing from 18 percent in 2003 to around 12 percent in 2006 and
2011. Given the strong correlation in the intergenerational transmission
of violent behavior, the decrease in this statistic may be informative
of the substantial drop in the incidence of violence observed in the
sample across survey waves. The differences in the distribution of
these pre-determined characteristics across waves reveal significant
changes in sample composition, and most importantly, on the types
of couples which remain in union.}\footnote{\textcolor{black}{See \citet{casique2014expresiones} for a rich analysis
of changes in household socio-economic characteristics and patterns
of intimate partner violence across the three survey waves.}}\textcolor{red}{{} }\textcolor{black}{Finally, levels of violence among
couples in the replication sample tend to be lower than those in the
pseudo-panel, consistent with the view that abuse in new couples may
be lower. To the extent that changes in the composition of beneficiary
households is driven by these marital selection dynamics, these sample
selection and treatment effect heterogeneity patterns help explain
the time path of spousal violence among beneficiary couples.}

\subsubsection*{\textcolor{black}{Alternative Explanations }}

\textcolor{black}{In the appendix, we evaluate and discuss other explanations
potentially consistent with the evidence: increasing rejection of
intimate partner violence, localized spillover effects among non-beneficiary
couples, changes in the de facto conditionality of the cash transfer
program, and generalized social violence mediating effects on spousal
abuse, among others.}

\section{Conclusion}

\label{sec:conc}

Our paper provides evidence of the relationship between the Oportunidades
CCT program, legal reforms codifying intimate partner violence as
legal grounds for unilateral divorce, and the prevalence of male-to-female
spousal violence in rural Mexico. More broadly, our paper contributes
to the debate on how the effects of cash transfer programs on IPV
are mediated by their local context, an understudied aspect of empirical
studies in the literature \citep{baranov2020theoretical}. Our findings
are consistent with the legal reforms both effectively increasing
women's threat points in a household bargaining framework, therefore
increasing the effectiveness of cash transfer policies in the short-run,
but perhaps most importantly, also enhancing the extent to which marital
selection shapes outcomes in the longer-run. In stark contrast to
the short-run relationships documented in \citet{bobonis2013public},
we find that, in the longer-run, women in beneficiary households are
as likely to experience physical and sexual abuse as non-beneficiary
women. We also document how the concurrent legal reforms that eased
women's ability to exit relationships led to higher divorce rates,
in particular for women who experienced violence in their marriage
and for women who were beneficiaries of the Oportunidades program.

Together, the results suggest that the policies reinforced each others'
capacity to reduce IPV. As further evidence of these policy complementarities,
our analysis revealed that the short-run negative relationship between
the CCT and IPV first documented in \citet{bobonis2013public} was
mostly concentrated in states that were early adopters of the legal
reforms. These complementarities appear to arise through changes in
behaviour in the short-run, but through marital selection in the longer
term. Our findings highlight that Oportunidades no longer appearing
to protect women against violence in the longer-run can thus be explained
by these violent relationships dissolving over time. We conclude that
the absence of a relationship between Oportunidades and IPV in
more recent surveys should therefore be viewed as a success rather
than as a failing of the policies. Finally, though our work is descriptive
by design and we cannot uncover the true causal impact of the interventions,
future work studying the relationship between cash transfers and IPV
should take into account the relevant legal environment, and the potential
for marital selection, when evaluating the long-run effects of these
programs.

\clearpage{}

%\bibliographystyle{aea}
%\bibliographystyle{aea}
%\bibliography{bcm_bib}
% Bibstyle aea.bst version 2009.05.20


% Bibstyle aea.bst version 2009.05.20

\begin{thebibliography}{xx}

\let\textbf\textrm

\harvarditem[Aguirre]{Edith Aguirre}{2019}{aguirre2019changes}
\textbf{Aguirre, Edith.} ``Do Changes in Divorce Legislation Have an Impact on
  Divorce Rates? The Case of Unilateral Divorce in Mexico.'' \emph{Latin
  American Economic Review} 28, no.~1 (2019):~9.

\harvarditem[Aizer]{Anna Aizer}{2010}{aizer2010gender}
\textbf{Aizer, Anna.} ``The Gender Wage Gap and Domestic Violence.''
  \emph{American Economic Review} 100, no.~4 (2010):~1847.

\harvarditem[Altonji, Elder and Taber]{Joseph~G Altonji, Todd~E Elder and
  Christopher~R Taber}{2005}{altonji2005selection}
\textbf{Altonji, Joseph~G, Todd~E Elder, {\rm and} Christopher~R Taber.}
  ``Selection on Observed and Unobserved Variables: Assessing the Effectiveness
  of Catholic Schools.'' \emph{Journal of Political Economy} 113, no.~1
  (2005):~151--184.

\harvarditem[Amaral, Bhalotra and Nishith]{Sofia Amaral, Sonia Bhalotra and
  Prakash Nishith}{2019}{amaral2019gender}
\textbf{Amaral, Sofia, Sonia Bhalotra, {\rm and} Prakash Nishith.} ``Gender,
  Crime and Punishment: Evidence from Women Police Stations in India.''
  \emph{Working Paper, University of Essex}  (2019).

\harvarditem[Anderberg et~al.]{Dan Anderberg, Helmut Rainer, Jonathan Wadsworth
  and Tanya Wilson}{2016}{anderberg2016unemployment}
\textbf{Anderberg, Dan, Helmut Rainer, Jonathan Wadsworth, {\rm and} Tanya
  Wilson.} ``Unemployment and Domestic Violence: Theory and Evidence.''
  \emph{The Economic Journal} 126, no.~597 (2016):~1947--1979.

\harvarditem[Anderson]{Siwan Anderson}{2018}{anderson2018legal}
\textbf{Anderson, Siwan.} ``Legal Origins and Female HIV.'' \emph{American
  Economic Review} 108, no.~6 (2018):~1407--39.

\harvarditem[Anderson and Genicot]{Siwan Anderson and Garance
  Genicot}{2015}{anderson2015suicide}
\textbf{Anderson, Siwan, {\rm and} Garance Genicot.} ``Suicide and Property
  Rights in India.'' \emph{Journal of Development Economics} 114
  (2015):~64--78.

\harvarditem[Angelucci and De~Giorgi]{Manuela Angelucci and Giacomo
  De~Giorgi}{2009}{angelucci2009indirect}
\textbf{Angelucci, Manuela, {\rm and} Giacomo De~Giorgi.} ``Indirect Effects of
  an Aid Program: How Do Cash Transfers Affect Ineligibles' Consumption?''
  \emph{American Economic Review} 99, no.~1 (2009):~486--508.

\harvarditem[Astorga and Shirk]{Luis Astorga and David~A
  Shirk}{2010}{astorga2010drug}
\textbf{Astorga, Luis, {\rm and} David~A Shirk.} ``Drug Trafficking
  Organizations and Counter-drug Strategies in the Us-Mexican Context.''
  \emph{Center for US-Mexican Studies Working Paper}  (2010).

\harvarditem[Avitabile]{Ciro Avitabile}{2012}{avitabile2012spillover}
\textbf{Avitabile, Ciro.} ``Spillover Effects in Healthcare Programs: Evidence
  on Social Norms and Information Sharing.'' \emph{Inter-American Development
  Bank Working Paper}  (2012).

\harvarditem[Baranov et~al.]{Victoria Baranov, Lisa Cameron, Diana
  Contreras~Suarez and Claire Thibout}{2020}{baranov2020theoretical}
\textbf{Baranov, Victoria, Lisa Cameron, Diana Contreras~Suarez, {\rm and}
  Claire Thibout.} ``Theoretical Underpinnings and Meta-analysis of the Effects
  of Cash Transfers on Intimate Partner Violence in Low-and Middle-Income
  Countries.'' \emph{The Journal of Development Studies}  (2020):~1--25.

\harvarditem[Beleche]{Trinidad Beleche}{2017}{beleche2017domestic}
\textbf{Beleche, Trinidad.} ``Domestic Violence Laws and Suicide in Mexico.''
  \emph{Review of Economics of the Household}  (2017):~1--20.

\harvarditem[Beleche and Lew]{Trinidad Beleche and Nellie
  Lew}{2011}{beleche2011til}
\textbf{Beleche, Trinidad, {\rm and} Nellie Lew.} ``Til Laws Do Us Part? The
  Impact of Changing Divorce Laws on Divorce Rates in Mexico.'' \emph{U.S. Food
  and Drug Administration, March. Unpublished Manuscript.}  (2011).

\harvarditem[BenYishay and Pearlman]{Ariel BenYishay and Sarah
  Pearlman}{2013}{benyishay2013homicide}
\textbf{BenYishay, Ariel, {\rm and} Sarah Pearlman.} ``Homicide and Work: The
  Impact of Mexico's Drug War on Labor Market Participation.''
  \emph{Unpublished manuscript, Vassar College}  (2013).

\harvarditem[Bloch and Rao]{Francis Bloch and Vijayendra
  Rao}{2002}{bloch2002terror}
\textbf{Bloch, Francis, {\rm and} Vijayendra Rao.} ``Terror as a Bargaining
  Instrument: A Case Study of Dowry Violence in Rural India.'' \emph{American
  Economic Review} 92, no.~4 (2002):~1029.

\harvarditem[Bobonis, Gonz{\'a}lez-Brenes and Castro]{Gustavo~J Bobonis,
  Melissa Gonz{\'a}lez-Brenes and Roberto Castro}{2013}{bobonis2013public}
\textbf{Bobonis, Gustavo~J, Melissa Gonz{\'a}lez-Brenes, {\rm and} Roberto
  Castro.} ``Public Transfers and Domestic Violence: The Roles of Private
  Information and Spousal Control.'' \emph{American Economic Journal: Economic
  Policy}  (2013):~179--205.

\harvarditem[Bobonis and Finan]{Gustavo~J Bobonis and Frederico
  Finan}{2009}{bobonis2009neighborhood}
\textbf{Bobonis, Gustavo~J, {\rm and} Frederico Finan.} ``Neighborhood Peer
  Effects in Secondary School Enrollment Decisions.'' \emph{The Review of
  Economics and Statistics} 91, no.~4 (2009):~695--716.

\harvarditem[Bobonis, Castro and Morales]{Gustavo~J Bobonis, Roberto Castro and
  Juan~S Morales}{2015}{bobonis2015conditionalc}
\textbf{Bobonis, Gustavo~J, Roberto Castro, {\rm and} Juan~S Morales.} 2015.
  ``Conditional Cash Transfers for Women and Spousal Violence: Evidence of the
  Long-term Relationship from the Oportunidades Program in Rural Mexico.'' IDB
  Working Paper Series N. 632.

\harvarditem[Bowlus and Seitz]{Audra~J Bowlus and Shannon
  Seitz}{2006}{bowlus2006domestic}
\textbf{Bowlus, Audra~J, {\rm and} Shannon Seitz.} ``Domestic Violence,
  Employment, and Divorce.'' \emph{International Economic Review} 47, no.~4
  (2006):~1113--1149.

\harvarditem[Brassiolo]{Pablo Brassiolo}{2016}{brassiolo2016domestic}
\textbf{Brassiolo, Pablo.} ``Domestic Violence and Divorce Law: When Divorce
  Threats Become Credible.'' \emph{Journal of Labor Economics} 34, no.~2
  (2016):~443--477.

\harvarditem[Buller et~al.]{Ana~Maria Buller, Amber Peterman, Meghna
  Ranganathan, Alexandra Bleile, Melissa Hidrobo and Lori
  Heise}{2018}{buller2018mixed}
\textbf{Buller, Ana~Maria, Amber Peterman, Meghna Ranganathan, Alexandra
  Bleile, Melissa Hidrobo, {\rm and} Lori Heise.} ``A Mixed-method Review of
  Cash Transfers and Intimate Partner Violence in Low-and Middle-income
  Countries.'' \emph{The World Bank Research Observer} 33, no.~2
  (2018):~218--258.

\harvarditem[Camacho]{Adriana Camacho}{2008}{camacho2008stress}
\textbf{Camacho, Adriana.} ``Stress and Birth Weight: Evidence from Terrorist
  Attacks.'' \emph{American Economic Review Papers and Proceedings} 98, no.~2
  (2008):~511--515.

\harvarditem[Casique and Castro]{Irene Casique and Roberto
  Castro}{2014}{casique2014expresiones}
\textbf{Casique, Irene, {\rm and} Roberto Castro.} ``Expresiones y contextos de
  la violencia contra las mujeres en M{\'e}xico: An\'{a}lisis comparativo de la
  ENDIREH 2011.'' \emph{Centro Regional de Investigaciones
  Multidisciplinarias-INMUJERES.}  (2014).

\harvarditem[Castro~Roberto and Medina]{Florinda~Riquer Castro~Roberto and
  Mar{\'\i}a~Eugenia Medina}{2006}{castro2006violencia}
\textbf{Castro~Roberto, Florinda~Riquer, {\rm and} Mar{\'\i}a~Eugenia Medina.}
  2006. \emph{Violencia de G{\'e}nero en las Parejas Mexicanas: resultados de
  la Encuesta nacional sobre la Din{\'a}mica de las relaciones en los Hogares
  2003.} Instituto Nacional de las Mujeres.

\harvarditem[Dee]{Thomas~S Dee}{2003}{dee2003until}
\textbf{Dee, Thomas~S.} ``Until Death Do You Part: The Effects of Unilateral
  Divorce on Spousal Homicides.'' \emph{Economic Inquiry} 41, no.~1
  (2003):~163--182.

\harvarditem[Dell]{Melissa Dell}{2015}{dell2015trafficking}
\textbf{Dell, Melissa.} ``Trafficking Networks and the Mexican Drug War.''
  \emph{American Economic Review} 105, no.~6 (2015):~1738--1779.

\harvarditem[Devries et~al.]{Karen~M Devries, Joelle~Y Mak, Loraine~J Bacchus,
  Jennifer~C Child, Gail Falder, Max Petzold, Jill Astbury and Charlotte~H
  Watts}{2013}{devries2013intimate}
\textbf{Devries, Karen~M, Joelle~Y Mak, Loraine~J Bacchus, Jennifer~C Child,
  Gail Falder, Max Petzold, Jill Astbury, {\rm and} Charlotte~H Watts.}
  ``Intimate Partner Violence and Incident Depressive Symptoms and Suicide
  Attempts: A Systematic Review of Longitudinal Studies.'' \emph{PLoS Medicine}
  10, no.~5 (2013):~e1001439.

\harvarditem[D\'{i}az and Saldarriaga]{Juan~Jos\'{e} D\'{i}az and Victor
  Saldarriaga}{2022}{diaz2022unconditional}
\textbf{D\'{i}az, Juan~Jos\'{e}, {\rm and} Victor Saldarriaga.}
  ``(Un)Conditional Love in the Time of Conditional Cash Transfers: The Effect
  of the Peruvian JUNTOS Program on Spousal Abuse.'' \emph{Economic Development
  and Cultural Change} 70, no.~2 (2022).

\harvarditem[Doyle and Aizer]{Joseph~J Doyle, Jr and Anna
  Aizer}{2018}{doyle2018economics}
\textbf{Doyle, Jr, Joseph~J, {\rm and} Anna Aizer.} ``Economics of Child
  Protection: Maltreatment, Foster Care, and Intimate Partner Violence.''
  \emph{Annual Review of Economics} 10 (2018):~87--108.

\harvarditem[Duflo]{Esther Duflo}{2012}{duflo2012women}
\textbf{Duflo, Esther.} ``Women Empowerment and Economic Development.''
  \emph{Journal of Economic Literature} 50, no.~4 (2012):~1051--1079.

\harvarditem[Fearon and Hoeffler]{James Fearon and Anke
  Hoeffler}{2014}{fearon2014benefits}
\textbf{Fearon, James, {\rm and} Anke Hoeffler.} ``Benefits and Costs of the
  Conflict and Violence Targets for the Post-2015 Development Agenda.''
  \emph{Conflict and Violence Assessment Paper, Copenhagen Consensus Center}
  (2014).

\harvarditem[Friedberg and Stern]{Leora Friedberg and Steven
  Stern}{2014}{friedberg2014marriage}
\textbf{Friedberg, Leora, {\rm and} Steven Stern.} ``Marriage, Divorce, and
  Asymmetric Information.'' \emph{International Economic Review} 55, no.~4
  (2014):~1155--1199.

\harvarditem[Garc\'{i}a-Ramos]{Aixa
  Garc\'{i}a-Ramos}{2021}{GARCIARAMOS2021102623}
\textbf{Garc\'{i}a-Ramos, Aixa.} ``Divorce Laws and Intimate Partner Violence:
  Evidence from Mexico.'' \emph{Journal of Development Economics}
  (2021):~102623.

\harvarditem[Heath, Hidrobo and Roy]{Rachel Heath, Melissa Hidrobo and Shalini
  Roy}{2020}{heath2020cash}
\textbf{Heath, Rachel, Melissa Hidrobo, {\rm and} Shalini Roy.} ``Cash
  transfers, polygamy, and intimate partner violence: Experimental evidence
  from Mali.'' \emph{Journal of Development Economics} 143 (2020):~102410.

\harvarditem[Heise]{Lori Heise}{2011}{heise2011works}
\textbf{Heise, Lori.} ``What Works to Prevent Partner Violence? An Evidence
  Overview.'' \emph{Working Paper, STRIVE Research Consortium, London School of
  Hygiene and Tropical Medicine}  (2011).

\harvarditem[Hoehn-Velasco and Penglase]{Lauren Hoehn-Velasco and Jacob
  Penglase}{2021}{hoehn2021impact}
\textbf{Hoehn-Velasco, Lauren, {\rm and} Jacob Penglase.} ``The Impact of
  No-fault Unilateral Divorce Laws on Divorce Rates in Mexico.'' \emph{Economic
  Development and Cultural Change} 70, no.~1 (2021):~203--236.

\harvarditem[Imbens]{Guido~W Imbens}{2003}{imbens2003sensitivity}
\textbf{Imbens, Guido~W.} ``Sensitivity to Exogeneity Assumptions in Program
  Evaluation.'' \emph{American Economic Review} 93, no.~2 (2003):~126--132.

\harvarditem[Imbens]{Guido~W Imbens}{2004}{imbens2004nonparametric}
\textbf{Imbens, Guido~W.} ``Nonparametric Estimation of Average Treatment
  Effects Under Exogeneity: A Review.'' \emph{Review of Economics and
  Statistics} 86, no.~1 (2004):~4--29.

\harvarditem[Kavanaugh, Sviatschi and Trako]{Guadalupe Kavanaugh, Maria
  Sviatschi and Iva Trako}{2019}{kavanaugh2019women}
\textbf{Kavanaugh, Guadalupe, Maria Sviatschi, {\rm and} Iva Trako.} ``Women
  Officers, Gender Violence and Human Capital: Evidence from Women's Justice
  Centers in Peru.'' \emph{Working Paper, Princeton University}  (2019).

\harvarditem[Lalive and Cattaneo]{Rafael Lalive and M~Alejandra
  Cattaneo}{2009}{lalive2009social}
\textbf{Lalive, Rafael, {\rm and} M~Alejandra Cattaneo.} ``Social Interactions
  and Schooling Decisions.'' \emph{The Review of Economics and Statistics} 91,
  no.~3 (2009):~457--477.

\harvarditem[Leon]{Gianmarco Leon}{2012}{leon2012civil}
\textbf{Leon, Gianmarco.} ``Civil Conflict and Human Capital Accumulation the
  Long-term Effects of Political Violence in Per{\'u}.'' \emph{Journal of Human
  Resources} 47, no.~4 (2012):~991--1022.

\harvarditem[Litwin, Perova and Reynolds]{Ashley Litwin, Elizaveta Perova and
  Sarah~Anne Reynolds}{2019}{litwin2019conditional}
\textbf{Litwin, Ashley, Elizaveta Perova, {\rm and} Sarah~Anne Reynolds.} ``A
  Conditional Cash transfer and Women's Empowerment: Does Bolsa Familia
  Influence Intimate Partner Violence?'' \emph{Social Science \& Medicine} 238
  (2019):~112462.

\harvarditem[Michaelsen and Salardi]{Maren~M Michaelsen and Paola
  Salardi}{2020}{michaelsen2015violence}
\textbf{Michaelsen, Maren~M, {\rm and} Paola Salardi.} ``Violence,
  Psychological Stress and Educational Performance During the ``War on Drugs''
  in Mexico.'' \emph{Journal of Development Economics} 143 (2020):~102387.

\harvarditem[Miguel and Kremer]{Edward Miguel and Michael
  Kremer}{2004}{miguel2004worms}
\textbf{Miguel, Edward, {\rm and} Michael Kremer.} ``Worms: Identifying Impacts
  on Education and Health in the Presence of Treatment Externalities.''
  \emph{Econometrica} 72, no.~1 (2004):~159--217.

\harvarditem[Miller and Segal]{Amalia~R Miller and Carmit
  Segal}{2019}{miller2019female}
\textbf{Miller, Amalia~R, {\rm and} Carmit Segal.} ``Do Female Officers Improve
  Law Enforcement Quality? Effects on Crime Reporting and Domestic Violence.''
  \emph{The Review of Economic Studies} 86, no.~5 (2019):~2220--2247.

\harvarditem[Perova]{Elizaveta Perova}{2010}{perova2010buying}
\textbf{Perova, Elizaveta.} ``Buying out of Abuse: How Changes in Women's
  Income Affect Domestic Violence.'' \emph{Unpublished manuscript, University
  of California-Berkeley}  (2010).

\harvarditem[Pierotti]{Rachael~S Pierotti}{2013}{pierotti2013increasing}
\textbf{Pierotti, Rachael~S.} ``Increasing Rejection of Intimate Partner
  Violence: Evidence of Global Cultural Diffusion.'' \emph{American
  Sociological Review} 78, no.~2 (2013):~240--265.

\harvarditem[Ritter~Burga]{Patricia Ritter~Burga}{2014}{ritter2014te}
\textbf{Ritter~Burga, Patricia.} ``M{\'a}s te quiero, m{\'a}s te pego? El
  efecto del Programa Juntos en el empoderamiento de las mujeres dentro del
  hogar.'' \emph{Working Paper. Instituto Nacional de Estad{\'\i}stica e
  Inform{\'a}tica. Lima}  (2014).

\harvarditem[Roy et~al.]{Shalini Roy, Melissa Hidrobo, John Hoddinott, Bastien
  Koch and Akhter Ahmed}{2019a}{roy2019can}
\textbf{Roy, Shalini, Melissa Hidrobo, John Hoddinott, Bastien Koch, {\rm and}
  Akhter Ahmed.} ``Can Transfers and Behavior Change Communication Reduce
  Intimate Partner Violence Four Years Post-program? Experimental Evidence from
  Bangladesh.'' \emph{IFPRI Discussion Paper 1869}  (2019).

\harvarditem[Roy et~al.]{Shalini Roy, Melissa Hidrobo, John Hoddinott and
  Akhter Ahmed}{2019b}{roy2019transfers}
\textbf{Roy, Shalini, Melissa Hidrobo, John Hoddinott, {\rm and} Akhter Ahmed.}
  ``Transfers, Behavior Change Communication, and Intimate Partner Violence:
  Postprogram Evidence from Rural Bangladesh.'' \emph{Review of Economics and
  Statistics} 101, no.~5 (2019):~865--877.

\harvarditem[Skoufias]{Emmanuel Skoufias}{2001}{skoufias2001progresa}
\textbf{Skoufias, Emmanuel.} ``PROGRESA and its Impacts on the Human Capital
  and Welfare of Households in Rural Mexico: A Synthesis of the Results of an
  Evaluation by IFPRI.'' \emph{International Food Policy Research Institute,
  Washington, DC}  (2001).

\harvarditem[Skoufias, Davis and De~La~Vega]{Emmanuel Skoufias, Benjamin Davis
  and Sergio De~La~Vega}{2001}{skoufias2001targeting}
\textbf{Skoufias, Emmanuel, Benjamin Davis, {\rm and} Sergio De~La~Vega.}
  ``Targeting the Poor in Mexico: An Evaluation of the Selection of Households
  into PROGRESA.'' \emph{World Development} 29, no.~10 (2001):~1769--1784.

\harvarditem[Stevenson and Wolfers]{Betsey Stevenson and Justin
  Wolfers}{2006}{stevenson2006bargaining}
\textbf{Stevenson, Betsey, {\rm and} Justin Wolfers.} ``Bargaining in the
  Shadow of the Law: Divorce Laws and Family Distress.'' \emph{The Quarterly
  Journal of Economics} 121, no.~1 (2006):~267--288.

\harvarditem[Tsaneva, Rockmore and Albohmood]{Magda Tsaneva, Marc Rockmore and
  Zahra Albohmood}{2018}{tsaneva2018effect}
\textbf{Tsaneva, Magda, Marc Rockmore, {\rm and} Zahra Albohmood.} ``The Effect
  of Violent Crime on Female Decision-making Within the Household: Evidence
  from the Mexican War on Drugs.'' \emph{Review of Economics of the Household}
  (2018):~1--32.

\harvarditem[Valdivia and Castro]{Marcos Valdivia and Roberto
  Castro}{2013}{valdivia2013gender}
\textbf{Valdivia, Marcos, {\rm and} Roberto Castro.} ``Gender Bias in the
  Convergence Dynamics of the Regional Homicide Rates in Mexico.''
  \emph{Applied Geography} 45 (2013):~280--291.

\end{thebibliography}


\pagebreak{}

\section*{\protect\pagebreak Tables and Figures}

\begin{figure}[H]
\begin{centering}
\caption{Date of introduction of legal reforms across Mexican states}
\par\end{centering}
\textcolor{white}{s}\\

\begin{centering}
\includegraphics[viewport=0bp 0bp 3600bp 2066bp,width=12cm]{Figures/timelineMap}
\par\end{centering}
\centering{}\textcolor{white}{s}\\
{\footnotesize{}Notes: Map shows the years in which the legal reforms
to state civil codes where introduced across Mexico.}{\footnotesize\par}
\end{figure}

\begin{figure}[H]
\begin{centering}
\caption{Introduction of legal reforms and divorce in the ENDIREH pseudo-panel}
\par\end{centering}
\begin{centering}
\includegraphics[width=16cm]{Figures/divrate_bl2}
\par\end{centering}
\centering{}{\small{}Notes: The figure shows the number of states
that had passed the law reform and the percentage of couples in the
ENDIREH 2011 pseudo-panel that had divorced between 1997 and 2008}{\small\par}
\end{figure}

\begin{figure}[H]
\begin{centering}
{\scriptsize{}\caption{Event study: Divorce laws and divorce rates for women in the ENDIREH
2011 pseudo-panel.}
}{\scriptsize\par}
\par\end{centering}
\begin{centering}
(a) All women
\par\end{centering}
\begin{centering}
{\scriptsize{}\includegraphics[width=12cm]{Figures/EVstudyDrate_weight}}{\scriptsize\par}
\par\end{centering}

\begin{centering}
(b) Beneficiaries of Oportunidades
\par\end{centering}
\begin{centering}
{\scriptsize{}\includegraphics[width=12cm]{Figures/EVstudyDBENEFrate_weight}}{\scriptsize\par}
\par\end{centering}
\begin{centering}
(c) Non-beneficiaries
\par\end{centering}
\begin{centering}
{\scriptsize{}\includegraphics[width=12cm]{Figures/EVstudyDNBENEFrate_weight}}{\scriptsize\par}
\par\end{centering}
{\scriptsize{}Notes: The figure shows the results of an event-study
analysis in which coefficients capture changes in divorce rates relative
to the year before the reform. Sample is a state-year panel of women
in the ENDIREH 2011 pseudo-panel using inverse sampling weights aggregated
at the state level. The first and last coefficients are binned to
capture all pre and post observations respectively. Top figure shows
all women, middle figure shows women who were beneficiaries of Oportunidades
and bottom panel shows women who were not beneficiaries of Oportunidades.}{\scriptsize\par}
\end{figure}

\begin{figure}[H]
\begin{centering}
\caption{Incidence of physical (top) and sexual violence (bottom) by beneficiary
status and survey wave; pseudo-panel (left) and replication sample
(right)}
\par\end{centering}
\begin{centering}
\includegraphics[width=8cm]{Figures/F1_violfi}\includegraphics[width=8cm]{Figures/F1_violfiRS}
\par\end{centering}
\begin{centering}
\includegraphics[width=8cm]{Figures/F1_violsexu}\includegraphics[width=8cm]{Figures/F1_violsexuRS}
\par\end{centering}
\centering{}%
\noindent\begin{minipage}[t]{1\columnwidth}%
{\footnotesize{}Notes: Shown are sample proportions of women reporting
being victims of physical violence (top) and sexual violence (bottom)
during the previous year, separately for beneficiaries and non-beneficiaries
of government social programs. The pseudo-panel includes couples in
rural villages with women ages 25/28/33 and older, with children aged
0-10/3-13/8-18, in a relationship since 1997. The replication sample
includes couples in rural villages with women ages 25 and older, with
children aged 0-10, who have been in a relationship for six years.
Sample proportions weighted by inverse sampling weights.}%
\end{minipage}
\end{figure}

\begin{figure}[H]
\begin{centering}
\caption{Acceptability of IPV and survey wave; pseudo-panel (left) and replication
sample (right)}
\par\end{centering}
\begin{centering}
\includegraphics[width=8cm]{Figures/F1_tieneDerecho}\includegraphics[width=8cm]{Figures/F1_tieneDerechoRS}
\par\end{centering}
\centering{}%
\noindent\begin{minipage}[t]{1\columnwidth}%
{\footnotesize{}Notes: Shown are sample proportions of women stating
that a husband ``has the right to hit his wife'', for each of the
2003/2006/2011 surveys. The pseudo-panel includes couples in rural
villages with women ages 25/28/33 and older, with children aged 0-10/3-13/8-18,
in a relationship since 1997. The replication sample includes couples
in rural villages with women ages 25 and older, with children aged
0-10, who have been in a relationship for six years. Sample proportions
weighted by inverse sampling weights. Sample proportions weighted
by inverse sampling weights.}%
\end{minipage}
\end{figure}

\pagebreak{}

\begin{table}[H]
\includegraphics[width=16.5cm]{Figures/t1_nn}

{\tiny{}Notes: Sample means weighted by inverse sampling weights.
The 2003 sample (column 1) includes couples with women ages 25 and
older, with children aged 0-10, and who have been in union since 1997
or earlier. The pseudo-panel (columns 2-3) includes couples with women
ages 28/33 and older, with children aged 3-13/8-18, and who have been
in union since 1997 or earlier, respectively for the 2006/2011 surveys.
The replication sample (columns 4-5) includes women ages 25 and older,
with children aged 0-10, and who have been in union for at least six
years.}{\tiny\par}
\end{table}

\begin{table}[H]
\begin{centering}
{\tiny{}\includegraphics[width=14cm]{Figures/t2_nn}}{\tiny\par}
\par\end{centering}
\centering{}{\tiny{}}%
\begin{minipage}[t]{0.85\columnwidth}%
{\tiny{}Notes: Sample means weighted by inverse sampling weights.
The 2003 sample (column 1) includes couples with women ages 25 and
older, with children aged 0-10, and who have been in union since 1997
or earlier. The pseudo-panel (columns 2-3) includes couples with women
ages 28/33 and older, with children aged 3-13/8-18, and who have been
in union since 1997 or earlier, respectively for the 2006/2011 surveys.
The replication sample (columns 4-5) includes women ages 25 and older,
with children aged 0-10, and who have been in union for at least six
years.}%
\end{minipage}{\tiny\par}
\end{table}

\begin{table}[H]
\begin{centering}
\includegraphics[width=16cm]{Figures/t3_nn}
\par\end{centering}
\centering{}{\tiny{}}%
\begin{minipage}[t]{0.95\columnwidth}%
{\tiny{}Notes: The exercise uses the state-panel constructed using
the ENDIREH 2011 survey for women in the pseudo-panel: women ages
33 and older, with children aged 8-18, and who were in union since
1997 or earlier. Other reforms controls include states' adoption of
the Penal Code Reform, which criminalized domestic violence, and the
Law of Access, Assistance and Prevention against Intra-Family Violence,
which established government assistance programs to prevent domestic
violence (Beleche, 2017). Regressions are weighted by sample size
at the state level. Robust standard errors in parentheses, clustered
at the state level; significant at ({*}) 90 percent, ({*}{*}) 95 percent,
({*}{*}{*}) 99 percent confidence levels.}%
\end{minipage}{\tiny\par}
\end{table}

\begin{table}[H]
\begin{centering}
\includegraphics[width=16cm]{Figures/t4_nn}
\par\end{centering}
\centering{}{\tiny{}}%
\begin{minipage}[t]{0.95\columnwidth}%
{\tiny{}Notes: Coefficient estimates from OLS regressions weighted
by survey sampling weights. Sample includes women both divorced and
in union for the ENDIREH 2011 pseudo-panel: women ages 33 and older,
with children aged 8-18, and who were in union since 1997 or earlier.
IPV Divorce Law (2003) is an indicator equal to 1 if the state had
passed the IPV Divorce Law by 2003. Individual controls include indicator
variables for woman and partner's age, woman and partner's indigenous
status, women's schooling-level indicators, the partner\textquoteright s
schooling attainment level, household size, cohabiting couple indicator,
years in union, and variables measuring reported histories of spousal
abuse in parental household during childhood. Robust standard errors
in parentheses, clustered at the state level; significant at ({*})
90 percent, ({*}{*}) 95 percent, ({*}{*}{*}) 99 percent confidence
levels. }%
\end{minipage}{\tiny\par}
\end{table}

\begin{table}[H]
\begin{centering}
\includegraphics[width=15cm]{Figures/t5_nn}
\par\end{centering}
\centering{}{\tiny{}}%
\begin{minipage}[t]{0.9\columnwidth}%
{\tiny{}Notes: Coefficient estimates from OLS regressions weighted
by survey sampling weights. The 2003 sample includes couples with
women ages 25 and older, with children aged 0-10, and who have been
in union since 1997 or earlier. The pseudo-panel includes couples
with women ages 28/33 and older, with children aged 3-13/8-18, and
who have been in union since 1997 or earlier, respectively for the
2006/2011 surveys. The replication sample includes women ages 25 and
older, with children aged 0-10, and who have been in union for at
least six years. IPV Divorce Law (2003) is an indicator equal to 1
if the state had passed the IPV Divorce Law by 2003. Individual controls
include indicator variables for woman and partner's age, woman and
partner's indigenous status, women's schooling-level indicators, the
partner\textquoteright s schooling attainment level, household size,
cohabiting couple indicator, years in union, and variables measuring
reported histories of spousal abuse in parental household during childhood.
Robust standard errors in parentheses, clustered at the state level;
significant at ({*}) 90 percent, ({*}{*}) 95 percent, ({*}{*}{*})
99 percent confidence levels. }%
\end{minipage}{\tiny\par}
\end{table}

\begin{table}[H]
\begin{centering}
\includegraphics[width=15cm]{Figures/t6_nn}
\par\end{centering}
\centering{}{\tiny{}}%
\begin{minipage}[t]{0.9\columnwidth}%
{\tiny{}Notes: Coefficient estimates from OLS regressions weighted
by survey sampling weights. The 2003 sample includes couples with
women ages 25 and older, with children aged 0-10, and who have been
in union since 1997 or earlier. The pseudo-panel includes couples
with women ages 28/33 and older, with children aged 3-13/8-18, and
who have been in union since 1997 or earlier, respectively for the
2006/2011 surveys. The replication sample includes women ages 25 and
older, with children aged 0-10, and who have been in union for at
least six years. IPV Divorce Law (2003) is an indicator equal to 1
if the state had passed the IPV Divorce Law by 2003. Individual controls
include indicator variables for woman and partner's age, woman and
partner's indigenous status, women's schooling-level indicators, the
partner\textquoteright s schooling attainment level, household size,
cohabiting couple indicator, years in union, and variables measuring
reported histories of spousal abuse in parental household during childhood.
Robust standard errors in parentheses, clustered at the state level;
significant at ({*}) 90 percent, ({*}{*}) 95 percent, ({*}{*}{*})
99 percent confidence levels.}%
\end{minipage}{\tiny\par}
\end{table}

\pagebreak{}


\end{document}

\section*{Appendix: For Online Publication}

\subsection*{Additional Information on the Oportunidades Program}

The Oportunidades program promotes children's human development in
education, nutrition, and health. The education component of Oportunidades
consists of subsidies typically provided to mothers, contingent on
their children's regular attendance at school.\footnote{Receipt of the education-specific benefits is contingent on children
attending school, which is verified by school personnel. For primary
and secondary school, the child becomes ineligible for support if
he or she misses school 4 times in a month without justification,
or 12 times during the school year. High school students become ineligible
if they are not certified as active during the school semester, defined
according to the regulations of the institution they attend.} Although PROGRESA initially targeted only children in primary and
middle school, Oportunidades was expanded to cover children in secondary
school. In 1998, these ranged from 70 to 255 pesos per month (approximately
7 to 25 USD), depending on the gender and grade level the child is
attending, with a maximum of 625 pesos (62.5 USD) per month per family.
Scholarship amounts have gradually increased, and in 2011 these ranged
from 150 pesos per month (approximately 12 USD), up to 960 pesos (77
USD).\footnote{This nominal average value of transfers has gradually increased since
the start of the program, and its purchasing power has varied (depending
on price levels in these areas and relative price changes with respect
to foreign currencies, i.e., USDs) throughout the 1997-2011 period.
Given these fluctuations, we opt to report the figure valid at the
date of the most recent ENDIREH survey, 2011.} Families also receive yearly benefits for the purchase of school
supplies of between 200 and 400 pesos (16 and 32 USD). In a further
expansion of the program in 2009, it now offers a cash transfer of
approximately 4,200 pesos to youth graduating from high school before
age 22 (J$ó$venes con Oportunidades).

The health and nutrition components consist of both cash transfers
and nutritional supplements. Supplements are targeted at infants 6-months
to 23-months old, pregnant and breast-feeding women, and children
aged 2-5 years who exhibit signs of malnutrition. Monthly cash transfers
for beneficiary families expanded throughout 1997-2011, by 2011 these
benefits included: nutritional support (Alimentario), 225 pesos (18
USD), originally part of PROGRESA; energy support (Energético), 60
pesos, established in 2007 to help families pay for energy costs (electricity,
gas, firewood, etc.); compensated nutritional support (Alimentario
Vivir Mejor), 120 pesos, established in 2008 to compensate families
for rising food prices; child support (Infantil Vivir Mejor), 105
pesos for every child aged between 0-9, established in 2010; elderly
support (Adultos Mayores), 315 pesos for every adult aged 70 or over,
established in 2006. These benefits are contingent on participation
by mothers in monthly health talks with the local health care provider,
the vaccination of family members, health checks of all children under
5 years old, and biannual health checks of all household members.
Overall, the program transfers are important, representing approximately
10 percent of the average expenditures of beneficiary families \citep{skoufias2001progresa}.
Maximum benefit levels have increased by approximately 20 percent
over time for families with children in only elementary or middle
school, but have almost doubled for those with children in secondary
school (see Figure 6).


\subsection*{Empirical Extensions and Alternative Explanations}

\textcolor{black}{We evaluate here a series of empirical extensions
and alternative explanations that can potentially help us understand
the empirical patterns we documented.}\textcolor{red}{}

\subsubsection*{Unobserved Heterogeneity Within-villages}

In order to address the potential concerns of unobserved heterogeneity
in the within-village household comparison, we pursue a set of tests
and sensitivity analyses inspired by the work on diagnostics of selection
on observable and unobservable variables \citep[i.e.,][]{imbens2003sensitivity,imbens2004nonparametric,altonji2005selection}.
Essentially, we identify which observable characteristics ($X_{iv}$)
that are correlated with treatment assignment ($T_{iv}$) - the woman's
age, partner's age, partner's schooling, family size, and years in
union - are also significant predictors of spousal abuse outcomes.
These may plausibly be the covariates most correlated with the unobservable
characteristics that jointly determine program eligibility/take-up
and violence outcomes. For those identified variables, we evaluate
the robustness of the results to flexible specifications that allow
for high-order and interaction terms between these variables, and
also include interactions with the woman's levels of education \citep{bobonis2013public}.
The results obtained from this sensitivity analysis are qualitatively
and quantitatively similar across specifications.

\subsubsection*{Increasing Rejection of Intimate Partner Violence}

\textcolor{black}{Recent research has documented a rapid global diffusion
of norms regarding the unacceptability of spousal violence across
a broad set of countries. Specifically, \citet{pierotti2013increasing}
uses nationally representative, repeated cross-sectional data from
Demographic and Health Surveys (DHS) across a broad set of low and
middle-income countries to document that women of reproductive age
have increasingly rejected the justification of violence from intimate
partners. She argues that new global cultural scripts rejecting violence
against women - via international and national policies and discussions
starting in the mid/late 1990s - may then be reflected in modifications
of individual attitudes towards IPV across a large spectrum of societies.
These new global scripts and norms may have also diffused across Mexican
society in such ways as to decrease women's tolerance for IPV. }\footnote{\textcolor{black}{Suggestive of this phenomenon in the Mexican context
is the passage of laws promoting gender equality and establishing
the right of women to live free of violence in 2006 and 2007, respectively.
Reports in the 2011 ENDIREH survey that 73 percent of women are knowledgeable
of the gender equality legislation and 82 percent of women report
being knowledgeable of the freedom from violence legislation are consistent
with a strong dissemination of these scripts as embodied in national
policy.}}

\textcolor{black}{To evaluate this hypothesis, we use additional information
available in the ENDIREH data. Following the analysis in \citet{pierotti2013increasing},
we construct an indicator variable that measures whether the woman
believes an intimate partner is justified in hitting or beating his
female partner when she does not meet her responsibilities.}\footnote{{The 2003 and 2006 survey rounds ask the same question:
\textquotedbl En su opini\'on, cuando la mujer no cumple con sus obligaciones,
el marido tiene el derecho de pegarle?\textquotedbl{} {[}In your opinion,
when a women does not meet her responsibilities, the partner has the
right to hit her?{]}. In contrast, the question in the 2011 survey
round is modified: \textquotedbl El hombre tiene el derecho de pegarle
a su esposa?\textquotedbl{} {[}Does a man have the right to hit his
partner?{]} Therefore, the responses in the 2011 survey round are
not strictly comparable to those in earlier rounds. We report these
in order to show a more complete picture, subject to this caveat.}}{{} These measures are imperfectly comparable to those
from existing DHS data.}\footnote{\textcolor{black}{Specifically, \citet{pierotti2013increasing} constructs
outcome variables derived from questions that asked respondents whether
it is okay for a man to hit or beat his wife under certain circumstances.
Specifically, the most common form of the question asked, \textquotedbl Sometimes
a husband is annoyed or angered by things which his wife does. In
your opinion, is a husband justified in hitting or beating his wife
in the following situations?\textquotedbl{} The five scenarios presented
to respondents were (1) if she goes out without telling him, (2) if
she neglects the children, (3) if she argues with him, (4) if she
refuses to have sex with him, and (5) if she burns the food.\textquotedbl{}
\citep[p. 248]{pierotti2013increasing}.}}\textcolor{black}{{} The proportion of women in our sample who reported
a husband ``has the right'' to hit his wife decreased from approximately
20 percent in 2003 to approximately 9 percent in 2006 and 3 percent
in 2011. Figure 6 shows the trend in this measure among couples in
the sample across the three survey years, by beneficiary status. Consistent
with the cross-country evidence, the proportion of women reporting
some justification of IPV shows a sharp reduction over this time period.
This stark change in the justification of IPV occurs among women in
both beneficiary and non-beneficiary households. For the pseudo-panel
we observe a decrease of 12.4 percentage points (54 percent; significant
at 95 percent confidence) - from 22.8 percent in 2003 to 10.4 percent
in 2006 - among women in beneficiary households, and a similar change
of 9.9 percentage points (58 percent; significant at the 95 percent
confidence level) among those in non-beneficiary households. We estimate
further decreases of proportional size between the 2006 and 2011 survey
rounds. These patterns also hold for the replication sample (Figure
6, right), as the forces driving these patterns are present across
Mexican society.}

Table 7 presents estimates of the relationship between Oportunidades
and Acceptability of IPV (analogous to those in Table 6 but with this
measure as the dependent variable). We find that the relationship
between the CCT program and acceptability of IPV is heterogenous across
Mexican states. In particular, the interaction between Oportunidades
and the IPV legal reforms is negative and significant, suggesting
that beneficiary women in reform states are 14.5 percentage points
less likely to report that husbands are justified in hitting their
wives, relative to those in non-reform states (column 2). This evidence
provides additional support to the complementary roles of legal reforms
and the CCT program in reducing IPV. The negative relationship persists
in 2006 but becomes statistically insignificant (column 4). In 2011,
we find an overall negative relationship between the Oportunidades
program and Acceptability of IPV for the replication sample (columns
5 and 6, Panel B).

\subsubsection*{\textcolor{black}{Spillover effects}}

\textcolor{black}{A first alternate explanation is that the program
empowers women in the community and provides them with the instruments
to prevent spousal abuse, directly via interactions with beneficiary
women with higher levels of empowerment in the community, or indirectly
via improved socio-economic conditions and options outside of current
relationships or changes in the norms of intolerance of abuse, among
other mechanisms.}\footnote{\textcolor{black}{As shown by \citet{angelucci2009indirect} and \citet{avitabile2012spillover},
the program had spillover effects on the consumption levels and health
behaviors (i.e., cervical cancer checks) of non-beneficiary households.
\citet{bobonis2009neighborhood}, \citet{lalive2009social} show evidence
of spillovers effects on middle school participation among children
in non-beneficiary households.}}\textcolor{black}{{} Therefore, to the extent that these spillover effects
reduce the incidence of abuse among non-beneficiary women and increase
female partners' intolerance of abuse, this can help explain the patterns
shown earlier.}

\textcolor{black}{We evaluate this alternate explanation empirically
by estimating empirical models that capture spillover effects at the
level of the village. Specifically, we estimate a variant of our main
empirical model (1). The regression equation incorporating these effects
is the following: 
\begin{equation}
Y_{ivm}=\theta_{1}T_{ivm}+\beta_{1}E[T_{-i,v,m}]+X_{ivm}\beta_{2}+\alpha_{m}+\varepsilon_{ivm}
\end{equation}
where the treatment indicator $T_{ivm}$ equals one for beneficiary
household $i$ in village $v$, municipality $m$ and is zero otherwise;
$E[T_{-i,v,m}]$ represents the proportion of beneficiary households
in the sample in village $v$ (excluding household $i$). This specification
incorporates the possibility that local spillovers are a (linear)
function of the proportion of beneficiary households in the village
\citep[e.g.,][]{miguel2004worms}. These potential effects are captured
by the $\beta_{1}$ term. We also estimate additional specifications
that allow for heterogeneous spillover effects among beneficiary and
non-beneficiary households by including an interaction term between
$T_{ivm}$ and $E[T_{-i,v,m}]$. Since the $E[T_{-i,v,m}]$ term is
highly collinear with village fixed effects, we substitute these for
municipality fixed effects in these specifications.}

\textcolor{black}{We report estimates of these models in the online
appendix Table 8, columns 1-3; Panel A reports estimates for the 2006
data whereas Panel B reports analogous ones for the 2011 data. The
estimates imply that a 10 percentage point increase in the proportion
of beneficiary women leads to a statistically insignificant 0.3 percentage
point increase in the incidence of physical or sexual abuse in 2006,
and a 0.2 percentage point decrease in its incidence in 2011 (column
1). In the specification allowing for heterogeneous spillover effects
by beneficiary status, the point estimates imply a statistically insignificant
decrease in spousal violence for women in villages with larger shares
of beneficiaries (column 3). Finally, it is worth noting that issues
of unobserved heterogeneity generally cause upward bias (in absolute
magnitude) in the estimates of the spillover effects; such that these
can be considered overestimates of the true spillover or social interaction
effects. We conclude that this alternative mechanism cannot explain
the results.}

\subsubsection*{\textcolor{black}{Generalized Social Violence}}

\textcolor{black}{Another potential alternate explanation we consider
is changes in the incidence of abuse or in its reporting due to the
marked increase in social violence. As is well known, Mexico has seen
a surge in homicide rates since 2007, and it has been concentrated
in particular regions of the country. Many analysts attribute this
drastic change in the level of violence to consequences of the federal
government's anti-crime policies meant to combat drug cartels \citep[e.g.,][]{astorga2010drug,dell2015trafficking}.}

\textcolor{black}{We consider the potential for this surge in generalized
social violence across municipalities and or states to affect the
trends in spousal abuse and the relationship with program beneficiary
status. On one hand, to the extent that the surge in homicides can
impinge on partners' stress levels or their levels of emotional health
more broadly, this could lead to greater conflict-related abuse. On
the other hand, if this generalized conflict negatively impinges on
women's willingness to report events of abuse, this would be consistent
with the significant drop in reported abuse rates in 2011, although
not so in the year 2006. These are two potential mechanisms that would
induce heterogeneity in the relationship, among others.}

\textcolor{black}{We evaluate this idea empirically by estimating
models that capture these factors at the level of the municipality
or state. The regression equation incorporating these factors is the
following: 
\begin{equation}
Y_{ivm}=\theta_{1}T_{ivm}+\theta_{2}T_{ivm}H_{m(s)}+\beta_{1}H_{m(s)}+X_{ivm}\beta_{2}+\varepsilon_{ivm}
\end{equation}
The $H_{m(s)}$ variable measures the homicide rate per hundred thousand
individuals in municipality $m$ (or alternatively, state $s$); the
other variables are defined as above.}\footnote{\textcolor{black}{The homicide data is available from Mexico's National
Statistics and Geography Institute (INEGI). We follow the standard
specifications in the literature and estimate the relationship between
violence and individual/household outcomes with measures of violence
at the municipality level \citep[e.g.,][]{camacho2008stress,leon2012civil}.
Empirical models of this sort find strong relationships with adult
labor force participation \citep{benyishay2013homicide} and student
achievement \citet{michaelsen2015violence} in Mexico.}}\textcolor{black}{{} The homicides measures are included for the calendar
year preceding the household survey (2005 and 2010, respectively)
since the surveys are conducted over a long time period and we aim
to ensure that the timing of measured homicides predates that of abuse
outcomes.}\footnote{\textcolor{black}{The results are robust to using contemporaneous
year measures (2006 and 2011, respectively). These are also robust
to the use of gender-specific homicide rates, in spite of the different
trends among the victim's gender, shown in \citet{valdivia2013gender}.
Estimates are available upon request.}}\textcolor{black}{{} The $\beta_{1}$ coefficient captures the partial
correlation between homicides and spousal abuse rates among non-beneficiary
couples, whereas the $\theta_{2}$ term captures the differential
correlation among beneficiary ones. In our main specification, because
the homicide rate is measured at the municipality level, we do not
include village fixed effects in this specification. In a second specification
with village fixed effects, we can identify the differential effect
for beneficiary couples.}

\textcolor{black}{We report estimates of these models fin Table 8,
columns 4-6. The estimate for 2006 implies that a one standard deviation
increase in the municipality-level homicide rate (10.18 deaths per
100,000 individuals) is associated with a 1.6 percentage point decrease
in spousal abuse (column 4, 0.01018$\times$-1.615), and although
the relationship is stronger for beneficiary women (column 6), this
is not differentially so. Analogous estimates for survey year 2011
imply that a one standard deviation increase in the homicide rate
is associated with a 0.5 percentage decrease in IPV, not statistically
significant (0.01971$\times-0.303$). We conclude that, although some
evidence is suggestive of generalized social violence being associated
with overall changes in IPV, and in particular for 2006, it does not
explain the loosening of the relationship between program beneficiary
status and spousal abuse.}\footnote{\textcolor{black}{These results are somewhat consistent with recent
work presented in \citet{tsaneva2018effect} that finds that the effects
of violent crime on female decision-making in Mexico are small and
tend to be short-lived. }}

\subsubsection*{\textcolor{black}{Improvement in Women's Labor Market Opportunities}}

\textcolor{black}{An extensive literature documents that increases
in a woman's relative labour opportunities, possibly by increasing
her bargaining power within the household by means of an improvement
in her outside option, can lead to lower levels of violence \citep[e.g.,][]{bowlus2006domestic,aizer2010gender,anderberg2016unemployment}.
If women's relative income-generating opportunities have improved
in Mexico over the last decade, this may help explain the strong decline
in the incidence of violence. To evaluate this hypothesis, we estimate
models analogous to \citet{aizer2010gender}'s that capture the mitigating
effects of the gender wage gap at the state level. The regression
equation incorporating these effects is the following: 
\begin{equation}
Y_{ivs}=\theta_{1}T_{ivs}+\theta_{2}T_{ivs}W_{s}+\beta_{1}W_{s}+X_{ivs}\beta_{2}+\varepsilon_{ivs}
\end{equation}
The $W_{s}$ variable measures the female/male wage ratio in state
$s$ relative to the average gender wage gap for the sample of women
and men in our study; the other variables are defined as above. The
$\beta_{1}$ coefficient captures the partial correlation between
the female/male wage ratio and spousal abuse rates among non-beneficiary
couples, whereas the $\theta_{2}$ term captures the differential
correlation among beneficiary ones. We use the state-level rural wage
gap measure because the surveys are representative at the state level
and thus the lowest level of aggregation at which these measures can
be consistently estimated is at this level.}\footnote{\textcolor{black}{Using a state-level female wage gap measure may
be somewhat restrictive for purposes of the analysis, as it may not
appropriately capture the relative labor market opportunities women
face across distinct municipalities and villages within the state.
However, it should capture broad differences at the state level in
these relative labor market opportunities.}}\textcolor{black}{{} Moreover, because the female/male wage ratio rate
is measured at the state level, we do not include village or state
fixed effects in this specification. In a second specification with
state fixed effects, we can identify the differential effect for beneficiary
couples.}

\textcolor{black}{We report estimates of these models in Table 8,
columns 7-9. The estimate for 2006 implies that a 10 percentage point
increase in the female/male wage ratio is associated with a decrease
in spousal abuse by 0.2 percentage points (Panel A, column 7). The
relationship for 2011 implies that an analogous increase in the later
period is associated with a 2.1 percentage point decrease in spousal
abuse (Panel B, column 7).}\footnote{\textcolor{black}{In 2006, the average wage gap for our sample was
24 percentage points, and by 2011, the average wage gap was substantially
smaller, at 12 percentage points.}}\textcolor{black}{{} Neither of these estimates is statistically significant.
Moreover, if we estimate a model that allows for a heterogeneous response
by couples' beneficiary status, it implies that beneficiary women
in states with lower female/wage ratios observe larger reductions
in spousal abuse, relative to beneficiary women in states with higher
female/male wage ratios. Though the estimates are not statistically
significant for 2006 (column 9, Panel A), they are significant at
the 90 percent level for 2011 (column 9, Panel B). The coefficients
imply, on one end of the spectrum, that beneficiary women in legal
reform states with the lowest female/male wage ratios (0.79), are
4.5 percentage points less likely to be victims of spousal abuse (0.79$\times0.319$
- 0.264 - 0.035). On the other hand, beneficiary women in non-reform
states with the highest female/male wage ratios (0.98), are 4.5 percentage
points more likely to be victims of spousal abuse (0.98$\times0.319$
- 0.264). These heterogeneous effects imply a greater reduction in
spousal abuse rates among beneficiary households in contexts of larger
gender wage gaps. It also suggests that the complementarity between
IPV divorce laws and the CCT program, though much weaker than in the
short-run, tends to persist in the longer-run.}

\pagebreak{}

\begin{figure}[H]
\caption{Oportunidades Program Maximum Monthly Benefits (in Real 1998 Mexican
Pesos)}

\includegraphics[width=16cm]{Figures/oport}

Source: Authors' calculations based on data from Secretariat of Social
Development (SEDESOL) Mexico and the Bank of Mexico.
\end{figure}

\begin{table}[H]
\includegraphics[width=15cm]{Figures/t7_nn}
\centering{}{\tiny{}}%
\begin{minipage}[t]{0.95\columnwidth}%
{\tiny{}Notes: Coefficient estimates from OLS regressions weighted
by survey sampling weights. Acceptability of IPV is an indicator on
whether the woman believes a partner is justified in hitting or beating
his female partner. The 2003 sample includes couples with women ages
25 and older, with children aged 0-10, and who have been in union
since 1997 or earlier. The pseudo-panel includes couples with women
ages 28/33 and older, with children aged 3-13/8-18, and who have been
in union since 1997 or earlier, respectively for the 2006/2011 surveys.
The replication sample includes women ages 25 and older, with children
aged 0-10, and who have been in union for at least six years. IPV
Divorce Law (2003) is an indicator equal to 1 if the state had passed
the IPV Divorce Law by 2003. Individual controls include indicator
variables for woman and partner's age, woman and partner's indigenous
status, women's schooling-level indicators, the partner\textquoteright s
schooling attainment level, household size, cohabiting couple indicator,
years in union, and variables measuring reported histories of spousal
abuse in parental household during childhood. Robust standard errors
in parentheses, clustered at the state level; significant at ({*})
90 percent, ({*}{*}) 95 percent, ({*}{*}{*}) 99 percent confidence
levels.}%
\end{minipage}{\tiny\par}
\end{table}

\begin{table}[H]
\begin{centering}
\includegraphics[width=16cm]{Figures/t8_nn}
\par\end{centering}
\centering{}{\tiny{}}%
\begin{minipage}[t]{0.95\columnwidth}%
{\tiny{}Notes: Coefficient estimates from OLS regressions weighted
by survey sampling weights. The 2003 sample includes couples with
women ages 25 and older, with children aged 0-10, and who have been
in union since 1997 or earlier. The pseudo-panel includes couples
with women ages 28/33 and older, with children aged 3-13/8-18, and
who have been in union since 1997 or earlier, respectively for the
2006/2011 surveys. The replication sample includes women ages 25 and
older, with children aged 0-10, and who have been in union for at
least six years. IPV Divorce Law (2003) is an indicator equal to 1
if the state had passed the IPV Divorce Law by 2003. Individual controls
include indicator variables for woman and partner's age, woman and
partner's indigenous status, women's schooling-level indicators, the
partner\textquoteright s schooling attainment level, household size,
cohabiting couple indicator, years in union, and variables measuring
reported histories of spousal abuse in parental household during childhood.
Robust standard errors in parentheses, clustered at the state level;
significant at ({*}) 90 percent, ({*}{*}) 95 percent, ({*}{*}{*})
99 percent confidence levels. }%
\end{minipage}{\tiny\par}
\end{table}

\pagebreak{}

\subsection*{Theoretical Framework}

We present a model of spousal relations to highlight the mechanisms
through which legal reforms and conditional cash transfers could affect
divorce and conflict within the household. Our theoretical framework
incorporates bargaining and conflict with incomplete information,
together with endogenous divorce decisions, to formalize the effects
that these two distinct sets of policies can have on the incidence
of conflict and divorce. For simplicity, we use a model where male
partners have significantly more bargaining power and this is fixed.
\footnote{An extension to increasing women\textquoteright s bargaining power
eliminates testable predictions and stresses even further the need
for empirical approaches.}

As is common in the literature on intra-household bargaining, partners
use their resources to generate a marital surplus and bargain over
its allocation; who owns the resources in the household matters by
affecting the spouses' outside options. Building on the work of \citet{anderson2015suicide}
and \citet{brassiolo2016domestic}, in order for bargaining to fail
some of the time we assume that spouses derive some private satisfaction
with the marriage, whose magnitude is unknown to their partner. We
further assume that when an offer is rejected, conflict ensues. This
comes at a cost to each spouse, and a cost whose magnitude is realized
only at the time of the conflict. The cost of conflict in our model
can be thought of as the physical and or psychological pain that each
spouse endures during an episode of violence. Its magnitude is uncertain
ex-ante since it depends on many factors, including each spouse's
ability to cope with it. Separation can only be achieved after going
through a period of marital conflict; after observing her private
cost of conflict, the wife may decide to end the relationship.

We model the legal reform as a reduction in the cost to divorce faced
by the wife, as these policies allow women to unilaterally initiate
divorce proceedings if they were victims of family violence. We also
model the introduction of the CCT as an increase in the wife's income,
consistent with the idea that these transfers were targeted towards
the female partner. The model therefore allows us to think about the
effects of these two policies and their interaction on the incidence
of conflict and divorce. Most importantly, the model allows us to
capture the selection mechanisms through which the policies may affect
the types of couples that remain married and those that choose to
divorce.

\subsection*{Preferences}

We assume three possible states in a relationship: cooperation, conflict,
and divorce. The utility functions of the partners depend on the status
of their marriage.

\textbf{Cooperation}

The utility function in the cooperation state depends on the share
of household resources each partner gets and their private gains to
marriage. Specifically, we assume that preferences for the wife and
husband are represented by the utility functions:
\begin{equation}
V^{w}(x,\theta_{w};I_{w})=u({\color{blue}{\color{black}x(I_{h}+I_{w})}})+\theta_{w}\ \&\ V^{h}(x,\theta_{h};I_{h})=u({\color{blue}{\color{black}(1-x)(I_{h}+I_{w})}})+\theta_{h},
\end{equation}
where $u(.)$ is a concave function representing the utility from
consumption, i.e. $u'(.)\geq0$ and $u"(.)\leq0$, and $\theta_{i},\ i\in\{w,h\}$
represent a personal satisfaction from marriage. Partners enjoy consumption
from their shared resources, where $I_{i},\ i\in\{w,h\}$ is the individual
income each partner contributes to the marriage and $x$ captures
the share of resources that are allocated to the wife's consumption
. We follow \citet{bloch2002terror,bobonis2013public,friedberg2014marriage}
in assuming that each spouse's personal utility from marriage is private
information; $\theta_{i},\ i\in\{w,h\}$ are independent and follow
a distribution $\theta_{w},\theta_{h}\sim G(.)$ with support $[\underline{\theta},\bar{\theta}]$,
such that $\underline{\theta}>0$, i.e. both husband and wife enjoy
positive utility from marriage.

\textbf{Conflict}

Households in the conflict state do not share resources; they enjoy
utility of consumption only from their own income, and incur a cost
of conflict $\kappa_{i},\ i\in\{w,h\}$; this captures both a private
cost of conflict as perceived by each partner as well as the intensity
of the conflict. For instance, a more violent husband will impose
a higher $\kappa_{w}$ on his wife. Specifically, the utility functions
in this state are represented by:

\begin{equation}
U^{w}(\theta_{w,}\kappa_{w};I_{w})=u(I_{w})+\theta_{w}-\kappa_{w}
\end{equation}
for the wife, and

\begin{equation}
U^{h}(\theta_{h},\kappa_{h};I_{h})=u(I_{h})+\theta_{h}-\kappa_{h}
\end{equation}
for the husband. We assume $\kappa_{w},\kappa_{h}\sim F(.)$ with
support $[\underline{\kappa},\bar{\kappa}]$, and $\underline{\kappa}\geq0$,
i.e. conflict will impose a loss in utility for both husband and wife.\footnote{One important distinction in our setup relative to that in \citet{anderson2015suicide}
is that couples in the conflict state continue to enjoy the private
gains to marriage.} 

\textbf{Divorce}

Partners who separate enjoy only their own income for consumption:

\begin{equation}
U_{D}^{w}(I_{w},\rho)=u(I_{w})+\rho
\end{equation}
for the wife, and

\begin{equation}
U_{D}^{h}(I_{h})=u(I_{h})
\end{equation}
 for the husband. In addition, the parameter $\rho$ captures the
additional utility or outside option from divorce for the wife such
that $\rho<0$ would represent a net loss from divorce. A reduction
in divorce costs are represented by an increase in $\rho$.

\subsection*{Timing}

Once married, each partner observes his/her level of private satisfaction
from the marriage. The husband then proposes a division of household
resources such that the wife gets share $x$ (and the husband gets
share $1-x$) of total household income. The wife then chooses to
accept or reject the husband's offer. If the wife accepts, they enter
the cooperative state and enjoy $V^{w}$ and $V^{h}$ as their utility.
If the wife rejects the offer, they enter the conflict state and learn
their costs $\kappa$. The wife can then choose to divorce the husband.
If she decides to stay married, they stay in the conflict state and
get $U^{w}$ and $U^{h}$, otherwise if they divorce, they get $U_{D}^{w}$
and $U_{D}^{h}$. For simplicity, we assume husbands choose $x$,
while wives choose whether to divorce or not.

\subsection*{Decisions}

We use backward induction to present the decisions at each stage:

\textbf{Wife's divorce decision:} The wife chooses to divorce if $U^{w}$<$U_{D}^{w}$,
or equivalently, the couple remains married if $\kappa_{w}\leq\theta_{w}-\rho$.

\textbf{Partners' cooperation or conflict decision:} If the wife rejects
an offer, her expected utility is given by: 

\begin{equation}
\begin{array}{cc}
E_{w}(\theta_{w};I_{w},\rho) & =F(\theta_{w}-\rho)[u(I_{w})+\theta_{w}]-\int_{-\infty}^{\theta_{w}-\rho}\kappa dF(\kappa)+[1-F(\theta_{w}-\rho)][u(I_{w})+\rho]\\
 & =u(I_{w})+\rho+F(\theta_{w}-\rho)(\theta_{w}-\rho)-\int_{-\infty}^{\theta_{w}-\rho}\kappa dF(\kappa)
\end{array}
\end{equation}
The expected utility captures both possible states after rejecting
an offer, conflict or divorce. A wife accepts the offer $x$ if:
\begin{equation}
V^{w}(x,\theta_{w};I_{w})\geq E_{w}(\theta_{w};I_{w},\rho)
\end{equation}

\textbf{Husband's offer decision:} Let $\tilde{\theta}(x;I_{w},\rho)$
be the value of $\theta_{w}$ such that $V^{w}(x,\theta_{w};I_{w})=E_{w}(\theta_{w};I_{w},\rho)$,
then a wife accepts an offer if and only if $\theta_{w}\geq\tilde{\theta}(x;I_{w},\rho)$,
where $\tilde{\theta}(x;I_{w},\rho)$ solves 
\begin{equation}
u({\color{blue}{\color{black}x(I_{h}+I_{w})}})+\theta_{w}=u(I_{w})+\rho+F(\theta_{w}-\rho)(\theta_{w}-\rho)-\int_{-\infty}^{\theta_{w}-\rho}\kappa dF(\kappa)
\end{equation}

Hence, $G[\tilde{\theta}(x;I_{w},\rho)]$ is the probability that
an offer $x$ is rejected from the husband's perspective. The husband\textquoteright s
expected utility if the wife rejects the offer is: 
\begin{equation}
E_{h}(\theta_{h};I_{h},\rho)=u(I_{h})+[\theta_{h}-E(\kappa)]\frac{1}{G[\tilde{\theta}(x;I_{w},\rho)]}\int_{-\infty}^{\tilde{\theta}(x;I_{w},\rho)}F(\theta_{w}-\rho)dG(\theta_{w})
\end{equation}

\begin{flushleft}
Anticipating the potential rejection of the offer by the wife, the
husband chooses an offer $x^{*}(\theta_{h};I_{h},I_{w},\rho)$ that
maximizes his expected utility:
\begin{equation}
\begin{array}{ccc}
h(x,\theta_{h};I_{h},I_{w},\rho) & = & (1-G[\tilde{\theta}(x;I_{w},\rho)])V^{h}(x,\theta_{h};I_{h}+I_{w})+G[\tilde{\theta}(x;I_{w},\rho)]E_{h}(\theta_{h};I_{h},\rho)\\
 & = & [1-G(\tilde{\theta})][u({\color{blue}{\color{black}(1-x)(I_{h}+I_{w})}})+\theta_{h}]+G(\tilde{\theta})u(I_{h})\\
 &  & +(\theta_{h}-E(\kappa))\int_{-\infty}^{\tilde{\theta}}F(\theta_{w}-\rho)dG(\theta_{w}).
\end{array}
\end{equation}
\par\end{flushleft}

Our main objects of interest are the \textbf{\textit{divorce rate}}
\textendash{} the proportion of partners who choose to divorce \textendash{}
and the \textbf{\textit{conflict rate}} \textendash{} those who remain
in the conflict state. The divorce rate ($D$) is given by the probability
that a wife\textcolor{red}{{} }\textcolor{black}{rejects an offer times
}the probability that she chooses to divorce in the case of rejection.
Accordingly, it is given by:

\begin{equation}
D(I_{h},I_{w},\rho)=\int_{-\infty}^{\infty}\int_{-\infty}^{\tilde{\theta}(x^{*}(\theta_{h};I_{h},I_{w},\rho);I_{w},\rho)}(1-F[\theta_{w}-\rho])dG(\theta_{w})dG(\theta_{h})
\end{equation}

The conflict rate is given by the probability that the wife\textcolor{red}{{}
}\textcolor{black}{rejects an offer }conditional on\textcolor{black}{{}
not divorcing}:
\begin{equation}
C(I_{h},I_{w},\rho)=\frac{\int_{-\infty}^{\infty}\int_{-\infty}^{\tilde{\theta}(x^{*}(\theta_{h};I_{h},I_{w},\rho);I_{w},\rho)}F[\theta_{w}-\rho]dG(\theta_{w})dG(\theta_{h})}{1-\int_{-\infty}^{\infty}\int_{-\infty}^{\tilde{\theta}(x^{*}(\theta_{h};I_{h},I_{w},\rho);I_{w},\rho)}(1-F[\theta_{w}-\rho])dG(\theta_{w})dG(\theta_{h})}.
\end{equation}


\subsection*{Predictions}

We are interested in studying how changes in the two parameters of
interest, a reduction in divorce costs due to the legal reform, $\rho$,
and the CCT program (represented by an increase in the wife's income),
affect the two outcomes of interest: the divorce rate ($D$) and the
conflict rate ($C$). For simplicity, we assume that once an offer
$x^{*}$ has been made by the husband, it remains fixed. We summarize
the main results here and present derivations and extensions in the
theoretical appendix.

\subsubsection*{Result 1:}

\textit{As divorce costs decrease, the equilibrium divorce rate increases
}($\frac{\partial D}{\partial\rho}>0$). There are two forces driving
this result: a direct effect in which the decrease in $\rho$ (i.e.
an increase in the outside option) makes couples who are in conflict
more likely to choose to divorce; and an indirect effect by which
the wife's expected utility from rejecting an offer becomes higher
and hence it is more likely for wives to choose to reject an offer
and enter the conflict or divorced state.

\subsubsection*{Result 2:}

\textit{The change in the equilibrium conflict rate with respect to
a decrease in divorce costs is ambiguous }$(\frac{\partial C}{\partial\rho}$
< or > 0). \textit{Women with sufficiently high costs of conflict
}($\kappa_{w}$) \textit{or sufficiently low private gains from marriage
}($\theta_{w}$) \textit{are more likely to exit the relationship}
($\frac{\partial C}{\partial\rho}$ $<0$ if $\kappa_{w}>\bar{\kappa_{w}}$
or $\theta_{w}<\text{\ensuremath{\underbar{\ensuremath{\theta_{w}}}}}$).
Again, there are two main forces driving this result. On one hand,
after the offer is rejected by the wife she is more likely to choose
to divorce instead of remaining in the conflict state. This effect
makes couples more likely to get divorced and less likely to be in
conflict. On the other hand, given a higher expected utility in the
conflict state ($E_{w}(\theta_{w};I_{w},\rho)$), the more likely
the wife is to reject her husband's offer, such that the proportion
of households in the cooperative state decreases. The overall magnitude
of the effect depends on the distribution of the wives' private gains
from marriage ($\theta_{w}$).

\textit{The types of couple remaining in the conflict state changes
as $\rho$ increases.} In particular, women with high costs of conflict
$\kappa_{w}$ and lower private gains from marriage $\theta_{w}$
are more likely to exit the relationship (as highlighted by the direct
effect) relative to women with lower $\kappa_{w}$ and higher $\theta_{w}$.
If conflict manifests itself in different forms across different levels
of $\kappa,$then this particular selection mechanism, through which
different types of couples remain in marriage after changes in divorce
laws, is an important channel through which the potential effects
arise.

\subsubsection*{Result 3:}

\textit{The effect of an increase in the wife's income on the equilibrium
divorce rate depends on the wife's ex ante outside option and her
partner's private gains to marriage. Specifically, there exist thresholds
$\rho_{H},$ $\rho_{L}$ and $\hat{\theta}$, such that: }
\begin{flushleft}
\begin{equation}
\frac{\partial D}{\partial I_{w}}\begin{cases}
\begin{array}{ccc}
>0 & \text{if} & \rho>\rho_{H}\\
<0 & \text{if} & \rho<\rho_{L}
\end{array}\end{cases}
\end{equation}
\par\end{flushleft}

How the wife's income affects divorce rates depends on the ex-ante
allocation of resources $x^{*}$ proposed by the husband and the utility
function from consumption $u$. Two important determinants of $x^{*}$
are the wife's initial outside option, which depends on $\rho$, and
the husband's private utility from the marriage $\theta_{h}$. Women
facing very low divorce costs ($\rho>\rho_{H}$) are compensated by
husbands through a high share $x^{*}$ of household resources. However,
since utility is diminishing in consumption, an increase in $I_{w}$
increases the utility of rejecting the husband's offer relative to
the utility of staying together, such that $\frac{\partial D}{\partial I_{w}}>0$.
On the other hand, if women face very high divorce costs ($\rho<\rho_{L}$),
the reverse is true.

For, $\rho_{L}<\rho<\rho_{H},$ changes in the divorce rate will depend
on the distribution of types $F$ and $G$.Husbands who enjoy a high
private utility from marriage $\theta_{h}$ offer favourable allocations
$x^{*}$ to their wives such that divorce is less likely for these
couples, and vice versa. Couples in which the husband enjoys greater
benefit from the marriage are more likely to remain together after
the wife receives the cash transfer. If husbands with higher gains
from marriage are also less likely to perpetrate violence, then this
highlights one important selection mechanism by which CCTs can affect
the composition of married couples.

\subsubsection*{Result 4:}

\textit{The effect of an increase in the wife's income on the equilibrium
conflict rate is ambiguous. The change in the conflict rate with respect
to the income of the wife }$\frac{\partial C}{\partial I_{w}}$ \textit{is
more negative for women with relatively higher costs of conflict $\kappa_{w}$
and for women with relatively lower utility of marriage $\theta_{w}.$}

Even though the overall effects of a CCT for women is ambiguous for
both conflict and divorce, this selection mechanism suggests that
different types of women will be in each of these states. As income
increases, women become more likely to reject the offers from their
husbands. Once women are in the conflict state, those facing high
costs of conflict choose to divorce, while those facing relatively
lower costs remain in conflict.

The results here follow from the same analysis as in the previous
point. Note in particular that the wife's income does not shift her
relative preferences between the conflict and divorce states. However,
changes in the wife's income affect the $types$ of couples that remain
in marriage, as given by their utility from marriage $\theta$ and
their costs of conflict $\kappa$.

With respect to the differential effect of female income for conflict
among women with relatively higher costs of conflict, imagine two
sets of women $H$ ($\kappa_{w}\geq\tilde{\kappa}$) and $L$ ($\kappa_{w}<\tilde{\kappa}$)
facing different costs of conflict stratified by the threshold level
$\tilde{\kappa}=\theta_{w}-\rho$. The conflict rate in the high cost
group decreases or increases at a slower rate than that for women
in the low costs group, as income changes. That is, $\frac{\partial C_{H}}{\partial I_{w}}<\frac{\partial C_{L}}{\partial I_{w}}$
. Note that the inverse is true for divorce, that is, $\frac{\partial D_{H}}{\partial I_{w}}>\frac{\partial D_{L}}{\partial I_{w}}$. 

\textit{The change in the conflict rate with respect to the income
of the wife }$\frac{\partial C}{\partial I_{w}}$ \textit{is smaller.
}As before, with two sets of women $H$ ($\theta_{w}\geq\tilde{\theta}$)
and $L$ ($\theta_{w}<\tilde{\theta}$) facing different utility of
marriage divided by $\tilde{\theta}<\kappa_{w}+\rho$, then conflict
in the low utility group decreases or increases at a slower rate than
that for women in the high utility group, as income changes. That
is, $\frac{\partial C_{L}}{\partial I_{w}}<\frac{\partial C_{H}}{\partial I_{w}}$
. Note that the inverse is true for divorce, that is, $\frac{\partial D_{L}}{\partial I_{w}}>\frac{\partial D_{H}}{\partial I_{w}}$.
As before, this selection mechanism suggests that different types
of women will be in each of these states. Once women are in the conflict
state, those facing low utility from marriage choose to divorce, while
those facing relatively higher utility choose to remain in the conflict
state.

This simple model of household relationships outlines the mechanisms
through which the policies of interest may affect conflict and divorce.
The model highlights some of the tradeoffs faced and suggests that
both legal reforms and CCTs lead to important selection mechanisms
which affect the $types$ of couples which remain in union. In particular,
couples in which the partners enjoy greater utility from marriage
$\theta$ and in which wives face lower costs from conflict $\kappa_{w}$
are more likely to stay together following the introduction of these
policies.
\end{document}
